
\documentclass[11pt]{article}

\usepackage[margin=1in]{geometry}
\usepackage[T1]{fontenc}
\usepackage[utf8]{inputenc}
\usepackage{lmodern}
\usepackage{microtype}
\usepackage{amsmath,amssymb,amsthm,mathtools}
\usepackage{array}
\usepackage{booktabs}
\usepackage{longtable}
\usepackage{enumitem}
\usepackage{xcolor}
\usepackage{hyperref}
\usepackage{listings}
\usepackage{stmaryrd}
\usepackage{tikz}
\usepackage{algorithm}
\usepackage{algpseudocode}
\usepackage{enumitem}

\hypersetup{colorlinks=true,linkcolor=blue!60!black,urlcolor=blue!60!black,citecolor=blue!60!black}

\lstdefinestyle{deppy}{
  basicstyle=\ttfamily\small,
  keywordstyle=\color{blue!60!black},
  commentstyle=\color{green!40!black},
  stringstyle=\color{red!50!black},
  showstringspaces=false,
  frame=single,
  breaklines=true,
  columns=fullflexible
}

\newcommand{\DepPy}{\textsc{DepPy}}
\newcommand{\TensorGuard}{\textsc{TensorGuard}}
\newcommand{\LiquidHaskell}{\textsc{LiquidHaskell}}
\newcommand{\Zthree}{\textsc{Z3}}
\newcommand{\Type}{\ensuremath{\mathsf{Type}}}
\newcommand{\Expr}{\ensuremath{\mathsf{Expr}}}
\newcommand{\Pred}{\ensuremath{\mathsf{Pred}}}
\newcommand{\Heap}{\ensuremath{\mathsf{Heap}}}
\newcommand{\Env}{\ensuremath{\mathsf{Env}}}
\newcommand{\Core}{\ensuremath{\mathsf{Core}}}
\newcommand{\Theory}{\ensuremath{\mathsf{Theory}}}
\newcommand{\cost}{\ensuremath{\mathsf{cost}}}
\newcommand{\gain}{\ensuremath{\mathsf{gain}}}
\newcommand{\trust}{\ensuremath{\mathsf{trust}}}
\newcommand{\support}{\ensuremath{\mathsf{support}}}
\newcommand{\py}[1]{\texttt{\detokenize{#1}}}
\newcommand{\sem}[1]{\llbracket #1 \rrbracket}
\newcommand{\liquid}[1]{\{\, #1 \,\}}
\newcommand{\Proved}{\ensuremath{\mathsf{Proved}}}
\newcommand{\Refuted}{\ensuremath{\mathsf{Refuted}}}
\newcommand{\Unknown}{\ensuremath{\mathsf{Unknown}}}
\newcommand{\Assumed}{\ensuremath{\mathsf{Assumed}}}
\newcommand{\Bounded}{\ensuremath{\mathsf{Bounded}}}

\theoremstyle{definition}
\newtheorem{definition}{Definition}[section]
\newtheorem{example}[definition]{Example}
\newtheorem{remark}[definition]{Remark}
\newtheorem{assumption}[definition]{Assumption}

\theoremstyle{plain}
\newtheorem{theorem}[definition]{Theorem}
\newtheorem{lemma}[definition]{Lemma}
\newtheorem{proposition}[definition]{Proposition}
\newtheorem{corollary}[definition]{Corollary}

\title{Dependent Typing as an Optimization Problem for Normal Python Code\\
\large Toward Maximal Automatic Signature Inference for Python, Tensors, and Library-Rich Programs in \DepPy}
\author{Standalone design monograph for the \DepPy{} effort}
\date{\today}

\begin{document}
\maketitle
\tableofcontents

\begin{abstract}
This document proposes a deliberately Pythonic account of dependent typing from
first principles. The central problem is not merely ``how do we check a user-
written dependent signature?'' but ``how do we automatically infer the richest
sound signature that ordinary Python code supports?'' The motivating target is a
surface result such as ``\py{a : Tensor | issorted(a)}'' and, more generally, the
ability to infer as many semantically valuable facts as possible about values,
functions, heap objects, protocols, and library calls while keeping the source
program visually close to normal Python. We therefore recast dependent typing as
an explicit optimization problem over candidate semantic facts: shape facts,
length facts, sortedness facts, permutation facts, protocol obligations, heap
invariants, equality witnesses, and theory-specific predicates should be
inferred, ranked, compressed, and verified rather than merely checked when
manually supplied.

The underlying claim is not that Python should be turned into a theorem prover
syntax, nor that every dynamic behavior should be pushed into a rigid
annotation discipline. Instead, we argue that dependent typing for Python
should be organized around six design commitments. First, the surface language
should remain as close as possible to ordinary Python, including classes,
protocols, exceptions, mutation, tensors, and idiomatic library code. Second,
annotations should be optional and sparse: the system should synthesize
contracts, qualifiers, witness candidates, shape relations, ordering facts,
protocol requirements, and heap invariants wherever possible. Third, the
internal type language should truly support full dependent structure---not just
refinement predicates but also dependent functions, dependent pairs, equality
types, indexed families, witness-carrying results, and theory-indexed library
contracts. Fourth, the verification engine should expose an explicit trusted
boundary in which \Zthree{} and related procedures discharge a decidable core,
while richer obligations remain first-class proof goals instead of being
silently erased. Fifth, the system should make optimization central: there are
always more candidate typing facts than can be profitably kept, so inference
must maximize useful type information under cost, explainability, trust, and
fragment constraints. Sixth, the predicate language and library-theory bank
should be continuously extensible without requiring a new LLM pass every time a
new theory is added; adding a theory should enrich the deterministic synthesis
and solver loop, not force regeneration of the whole program analysis pipeline.

The emphasis of this monograph is therefore slightly different from a purely
language-theoretic treatment. We still define syntax, judgments, and semantic
objects carefully, but we place equal weight on extraction, ranking, library
modeling, protocol reasoning, contract compression, theory accretion, and novel
uses of \Zthree{} for maximal signature inference. We draw inspiration from
\LiquidHaskell{} for qualifier generation and SMT-backed checking, but we push
the idea outward to a wider design space: nominal classes, structural
protocols, heap facts, tensor theories such as PyTorch, effect summaries,
proof-carrying dependent specifications, and open-world theory packs that can be
added incrementally without rerunning any LLM. The target user is a Python
programmer or LLM writing Python-looking code, not a user expected to rewrite
their program into a proof assistant dialect.
\end{abstract}

\section{Research questions, thesis claims, and concrete contributions}

To make this document function more like a thesis than a position paper, it is
useful to state the central research questions explicitly.

\subsection{Research questions}
\begin{enumerate}[label=\textbf{RQ\arabic*.}]
  \item Can dependent typing for a language as dynamic and library-driven as
  Python be reformulated so that the primary output of analysis is not a single
  principal type, but a frontier of maximally informative semantic signatures?
  \item Can ordinary unannotated Python code---especially PyTorch-heavy
  code---support the inference of semantic facts such as sortedness,
  permutation, shape, protocol conformance, heap consistency, and witness-
  carrying results without forcing programmers into a proof-assistant surface
  syntax?
  \item Can the addition of new predicates and new library theories be made
  monotone and deterministic, so that theory growth improves analysis quality
  without requiring another LLM pass over the original program?
  \item Can \Zthree{} be used not merely as a checker for downstream
  obligations, but as the central optimizer for signature packing, frontier
  extraction, theory gating, and counterexample-guided theory growth?
  \item Can such a system approach the semantic ambition of full dependent
  types while remaining honest about the trusted fragment, surfacing residual
  proof obligations whenever automation stops?
\end{enumerate}

\subsection{Thesis claims}
The core thesis of this monograph is the conjunction of four claims.
\begin{enumerate}[label=\arabic*.]
  \item \textbf{Frontier claim.} The right semantic target for practical
  dependent typing in Python is not a singleton type but a frontier of
  candidate signatures ordered by informativeness, proof burden, and
  presentation cost.
  \item \textbf{Optimization claim.} The central inference problem is best
  expressed as an optimization problem over semantic atoms, relations, witness
  schemas, and theory activations rather than as pure derivation in a fixed
  syntax-directed system.
  \item \textbf{Accretion claim.} Theories should accumulate monotonically:
  every new theory pack enlarges the candidate basis and can only improve the
  achievable frontier, provided the installed pack is sound with respect to its
  declared trust boundary.
  \item \textbf{Pythonic claim.} Full dependent ambition and normal Pythonic
  surface syntax are compatible if the system treats annotations as optional
  views of an underlying inferred semantic object rather than as the sole source
  of truth.
\end{enumerate}

\subsection{Concrete contributions claimed by this document}
Relative to prior practical SMT-backed refinement systems, the document advances
several concrete contributions.
\begin{enumerate}[label=\textbf{C\arabic*.}]
  \item A \emph{frontier-valued typing judgment} in which the result of
  analysis is a set of undominated semantic signatures rather than a single
  derivation artifact.
  \item A unified \emph{maximal signature packing} objective that optimizes
  jointly over structural facts, semantic relations, witness schemas, and
  theory activation decisions.
  \item An \emph{open-world theory pack} model that makes library growth
  additive and deterministic rather than prompt-driven.
  \item A family of \emph{novel \Zthree{} uses} for lexicographic signature
  saturation, frontier enumeration, anti-chain extraction, hitting-set-guided
  theory growth, and theory gating.
  \item A \emph{pythonic operational stance} in which ordinary programs are
  lifted into semantic signatures first and only then optionally compressed back
  into annotations, decorators, or proof-goal payloads.
  \item A \emph{clean-slate implementation agenda} for rebuilding the
  infrastructure around these ideas rather than merely extending a refinement
  checker with more ad hoc heuristics.
\end{enumerate}

\subsection{What success would look like}
The document will be successful if a future implementation can take ordinary
Python such as
\begin{lstlisting}[style=deppy,language=Python]
def f(a):
    values, _ = torch.sort(a)
    return values
\end{lstlisting}
and infer not merely ``returns a tensor'' but a frontier containing objects as
strong as
\[
(a : \mathsf{Tensor}(n,\delta,d)) \to
(result : \mathsf{SortedTensor}(n,\delta,d) \mid
\mathsf{sorted}(result) \wedge \mathsf{perm}(result,a)).
\]
The essential thesis-like point is that this object should be regarded as the
primary semantic result, while any compact user-facing decorator or signature is
only a later rendering of it.


\section{Design stance: dependent types should feel native to Python}

The central design thesis of this file is that dependent typing for Python
should optimize for \emph{ordinary code first}. A successful system should be
able to look at scripts that use straightforward function definitions, idiomatic
conditionals, classes with \py{__init__}, data-carrying objects, exception-based
control flow, and familiar library APIs. It should then recover precise types
and contracts without first demanding that the user become an expert in a new
annotation metalanguage. This is a stricter requirement than mere convenience.
It changes the architecture of the entire system. If the surface language must
remain visually close to Python, then the source of typing information cannot be
only explicit annotations. It must also come from guards, branches, protocol
uses, index operations, attribute writes, exception exits, library calls,
container destructuring, and local proof obligations generated from dangerous
operations.

A useful way to summarize the goal is the following slogan:
\begin{quote}
Dependent typing for Python should infer as much as possible, require as little
as possible, and expose exactly what remains unproved.
\end{quote}
This slogan implies that the boundary between ``type checking'' and ``program
analysis'' must be softened. In a language like Idris, the user often writes the
dependent structure explicitly; in Python, the system must discover a large part
of it from the program text and from library-specific semantic theories. The
relevant comparison is therefore not only to classical dependent type systems,
but also to \LiquidHaskell{}, abstract interpretation, contract inference,
protocol checking, and solver-guided program analysis.

\begin{definition}[Pythonic dependent typing]
A \emph{pythonic dependent typing system} is a system whose typing judgments are
defined over ordinary Python surface syntax, whose principal information sources
include both explicit annotations and automatically harvested program facts, and
whose user-visible output may be any mixture of inferred contracts, suggested
minimal annotations, proof obligations, counterexamples, and verified core
verdicts.
\end{definition}

\begin{remark}
This definition is intentionally broader than ``refinement typing'' alone.
Refinement types are one important fragment, but the system also needs dependent
functions, dependent pairs, equality witnesses, indexed structures, and theory-
specific predicates supplied by library semantics.
\end{remark}



\section{Surface language priorities: normal Python, optional annotations, explicit leftovers}

We insist on three practical priorities.

\subsection{Priority 1: code should stay visually normal}
The preferred user script should look like a normal module. Consider the
following style of code:
\begin{lstlisting}[style=deppy,language=Python]
class Box:
    value: int

    def __init__(self, value: int) -> None:
        self.value = value

    def get(self):
        return self.value

def first(xs):
    return xs[0]
\end{lstlisting}
The first function uses one ordinary class annotation and one omitted method
annotation. The second function uses no annotation at all. A useful dependent
system should still discover that \py{first} needs a non-emptiness precondition,
that \py{Box.get} returns an \py{int}, and that writing \py{self.value = "oops"}
would violate the inferred or declared heap invariant. None of these require the
user to abandon Pythonic style.

\subsection{Priority 2: annotations are hints, not a precondition for precision}
The role of annotations changes in this design. They are not the only source of
truth; they are high-value candidate facts in an optimization landscape. If the
user supplies a return annotation, the system should exploit it. If the user
supplies a dependent signature string, the system should check and preserve it.
But if the user supplies nothing, the system should still synthesize useful
information from guards, uses, library theories, and generated proof
obligations.

\subsection{Priority 3: whatever cannot be proved must not disappear}
A common failure mode in practical verification systems is to silently erase the
most interesting semantic claims because they are hard to automate. We reject
that design. If a property such as permutation preservation for sorting or exact
heap ownership cannot be fully discharged in the current core, it should remain
visible as a proof goal. This is how the system remains faithful to full
dependent expressivity without pretending that the whole of Python semantics has
already been reduced to a trivial SMT fragment.



\section{From refinements to full dependent structure}

A recurring misunderstanding in work on practical dependent typing for dynamic
languages is that ``dependent typing'' simply means ``add predicates to base
types.'' That fragment is valuable but incomplete. The real target language of
\DepPy{} must support at least the following constructors.

\begin{enumerate}[label=\arabic*.]
  \item \textbf{Refinement types}: \(\liquid{x : \tau \mid \phi(x)}\).
  \item \textbf{Dependent functions}: \((x : A) \to B(x)\).
  \item \textbf{Dependent pairs}: \((x : A) \times B(x)\), exposed operationally as
  witness-carrying return values.
  \item \textbf{Equality types}: \(\mathsf{Eq}(A, u, v)\), needed for relating library
  outputs to inputs and for proof-carrying rewrites.
  \item \textbf{Indexed families}: tensors indexed by shape, vectors indexed by
  length, finite sets indexed by a bound, heaps indexed by permissions or
  regions, protocol environments indexed by method availability, and so on.
  \item \textbf{Theory-indexed propositions}: domain-specific predicates for
  PyTorch shapes, devices, sorting correctness, batching invariants, layout
  relations, database schemas, protocol conformance, and effect summaries.
\end{enumerate}

\begin{definition}[Dependent information object]
A \emph{dependent information object} is any typed artifact inferred or supplied
by the system whose semantics depends on values. This includes refinements,
indexed result types, witness-returning Sigma types, equality claims, and
library-theoretic contracts.
\end{definition}

\begin{example}[Non-refinement dependent result]
A Python function that exposes a witness can be modeled by a dependent pair:
\[
\mathsf{expose\_length\_and\_values} :
(n : \mathsf{Nat}) \to \mathsf{Vec}(\mathsf{Int}, n)
\to (k : \mathsf{Nat}) \times \mathsf{Vec}(\mathsf{Int}, k).
\]
Even if the current implementation presents this using a compact Python-side
signature string, the semantic object is genuinely beyond mere refinement
predicates.
\end{example}

\begin{remark}
The system therefore cannot be described as ``only Liquid Types for Python.'' It
must contain a Liquid-inspired inference core embedded inside a larger dependent
language.
\end{remark}



\section{Why typing should be formulated as an optimization problem}

In a practical Python program, the system usually has access to many more
candidate facts than it can profitably keep. A function body, its call sites,
its library models, and its control-flow structure may jointly imply dozens or
hundreds of plausible predicates, structural constraints, protocol facts,
witness candidates, and heap relations. Retaining all of them naively may be
logically sound but pragmatically bad: it increases solver cost, inflates
explanations, harms minimal annotation synthesis, and may push obligations into
harder fragments than the user actually needs.

This means that the core inference problem is not only:
\begin{quote}
Which typing facts are derivable?
\end{quote}
but rather:
\begin{quote}
Which derivable typing facts maximize useful precision under trust, explainability,
annotation burden, and solver-fragment constraints?
\end{quote}

\begin{definition}[Global typing selection problem]
Let $C$ be the finite set of candidate dependent information objects induced by a
program. Each candidate $c \in C$ carries metadata
\[
\mathsf{meta}(c) = \langle
\mathsf{scope}(c),
\mathsf{origin}(c),
\mathsf{support}(c),
\mathsf{fragment}(c),
\mathsf{burden}(c),
\mathsf{stability}(c)
\rangle.
\]
The \emph{global typing selection problem} is to choose $S \subseteq C$ so as to
maximize useful type information and minimize proof burden.
\end{definition}

A generic objective is:
\[
\max_{S \subseteq C}
\Big(
\lambda_1 \gain_{\mathrm{precision}}(S)
+
\lambda_2 \gain_{\mathrm{diagnostic}}(S)
+
\lambda_3 \gain_{\mathrm{reuse}}(S)
-
\lambda_4 \cost_{\mathrm{solver}}(S)
-
\lambda_5 \cost_{\mathrm{annotation}}(S)
-
\lambda_6 \cost_{\mathrm{instability}}(S)
\Big)
\]
subject to soundness and trusted-core constraints.

This objective is deliberately broad. It includes not only precision but also
reuse and human burden. A contract that is slightly weaker yet explainable,
portable, and solver-friendly may be preferable to a maximally strong but opaque
one. The system's job is to make that tradeoff explicit rather than accidental.



\section{Liquid-Haskell-style ideas, exploded outward}

The closest influential predecessor for the inference side of this design is
\LiquidHaskell{}. Its key ideas are powerful: infer qualifiers from program
syntax, generate implication constraints, solve them using an SMT engine, and
obtain automated refinement typing without requiring the user to write every
predicate by hand. We adopt this spirit, but we deliberately broaden its role.

\subsection{Classical Liquid-style components}
A standard Liquid-style pipeline provides:
\begin{enumerate}[label=\arabic*.]
  \item qualifier templates,
  \item subtyping constraints,
  \item Horn-like implication checks,
  \item SMT-based discharge,
  \item predicate selection guided by validity.
\end{enumerate}
These pieces remain valuable in Python and should be retained.

\subsection{Where Python forces a larger design}
However, Python adds realities that classical Liquid-style settings can often
ignore or contain more tightly:
\begin{itemize}
  \item open-world libraries,
  \item mutation and aliasing,
  \item nominal classes and structural protocols,
  \item effectful methods and heap updates,
  \item gradual user intent,
  \item library-specific theories such as tensor shape/device semantics,
  \item highly sparse annotations,
  \item LLM-authored but machine-checked guarantee payloads.
\end{itemize}
Therefore the Liquid core should not be the whole system. It should be a
component inside a larger optimization-and-theory architecture.

\subsection{Explosion points}
We ``explode'' the Liquid idea along several axes.

\paragraph{Explosion 1: from refinements to all dependent artifacts.}
Qualifiers are no longer only predicates over scalar variables. They may help
choose witness-returning Sigma types, theory-indexed library contracts,
protocol requirements, and heap summaries.

\paragraph{Explosion 2: from local expressions to library theories.}
A PyTorch call is not just a function application. It belongs to a semantic
theory containing shape equations, device constraints, permutation properties,
and API-level postconditions.

\paragraph{Explosion 3: from solving constraints to choosing what to keep.}
Derivability is not the end of inference. We still need to select which facts
become user-facing contracts, which remain internal, and which are retained only
as proof obligations.

\paragraph{Explosion 4: from programmer-written annotations to mixed evidence.}
The candidates may come from explicit decorators, harvested guards, data-flow
obligations, library signatures, protocol uses, heap accesses, or LLM-generated
formal guarantee payloads.


\section{Qualifier generation for sparse and zero-annotation programs}

If annotations are optional, candidate generation must be aggressive. A good
qualifier language for Python should still be finite and solver-friendly, but it
should be rich enough to cover ordinary idioms. The following families are
particularly important.
\subsection{Qualifier family: array bounds}
The system introduces qualifiers corresponding to \(0 \le i < \mathrm{len}(xs)\) from indexing, slicing, and update operations. In a standard
Liquid-style design, such qualifiers would be mined from syntax and then used
to instantiate implication constraints. In our expanded setting, they are also
fed into the global optimization layer, where they compete with richer
dependent artifacts such as protocol summaries, witness-returning signatures,
and library-theory lemmas. This is especially important in Python, because a
user may write no annotations at all, yet the code still contains local uses
that force semantic obligations. The system should therefore create the
candidate even when it has not yet decided whether to surface it as a user
precondition, a hidden proof obligation, or a theorem to be delegated to a
library theory solver.

\subsection{Qualifier family: nullability}
The system introduces qualifiers corresponding to \(x \neq \mathsf{None}\) from dereferences, method calls, and attribute reads. In a standard
Liquid-style design, such qualifiers would be mined from syntax and then used
to instantiate implication constraints. In our expanded setting, they are also
fed into the global optimization layer, where they compete with richer
dependent artifacts such as protocol summaries, witness-returning signatures,
and library-theory lemmas. This is especially important in Python, because a
user may write no annotations at all, yet the code still contains local uses
that force semantic obligations. The system should therefore create the
candidate even when it has not yet decided whether to surface it as a user
precondition, a hidden proof obligation, or a theorem to be delegated to a
library theory solver.

\subsection{Qualifier family: non-emptiness}
The system introduces qualifiers corresponding to \(\mathrm{len}(xs) > 0\) from head-like operations, iteration witnesses, and destructuring. In a standard
Liquid-style design, such qualifiers would be mined from syntax and then used
to instantiate implication constraints. In our expanded setting, they are also
fed into the global optimization layer, where they compete with richer
dependent artifacts such as protocol summaries, witness-returning signatures,
and library-theory lemmas. This is especially important in Python, because a
user may write no annotations at all, yet the code still contains local uses
that force semantic obligations. The system should therefore create the
candidate even when it has not yet decided whether to surface it as a user
precondition, a hidden proof obligation, or a theorem to be delegated to a
library theory solver.

\subsection{Qualifier family: protocol availability}
The system introduces qualifiers corresponding to \(\mathsf{hasMethod}(x,m)\) and \(\mathsf{hasField}(x,f)\) from attribute use. In a standard
Liquid-style design, such qualifiers would be mined from syntax and then used
to instantiate implication constraints. In our expanded setting, they are also
fed into the global optimization layer, where they compete with richer
dependent artifacts such as protocol summaries, witness-returning signatures,
and library-theory lemmas. This is especially important in Python, because a
user may write no annotations at all, yet the code still contains local uses
that force semantic obligations. The system should therefore create the
candidate even when it has not yet decided whether to surface it as a user
precondition, a hidden proof obligation, or a theorem to be delegated to a
library theory solver.

\subsection{Qualifier family: heap field consistency}
The system introduces qualifiers corresponding to \(\mathsf{fieldType}(o,f)=\tau\) from class declarations and writes to \py{self.f}. In a standard
Liquid-style design, such qualifiers would be mined from syntax and then used
to instantiate implication constraints. In our expanded setting, they are also
fed into the global optimization layer, where they compete with richer
dependent artifacts such as protocol summaries, witness-returning signatures,
and library-theory lemmas. This is especially important in Python, because a
user may write no annotations at all, yet the code still contains local uses
that force semantic obligations. The system should therefore create the
candidate even when it has not yet decided whether to surface it as a user
precondition, a hidden proof obligation, or a theorem to be delegated to a
library theory solver.

\subsection{Qualifier family: tensor shapes}
The system introduces qualifiers corresponding to shape equalities, broadcasting constraints, matrix compatibility, flatten/view obligations. In a standard
Liquid-style design, such qualifiers would be mined from syntax and then used
to instantiate implication constraints. In our expanded setting, they are also
fed into the global optimization layer, where they compete with richer
dependent artifacts such as protocol summaries, witness-returning signatures,
and library-theory lemmas. This is especially important in Python, because a
user may write no annotations at all, yet the code still contains local uses
that force semantic obligations. The system should therefore create the
candidate even when it has not yet decided whether to surface it as a user
precondition, a hidden proof obligation, or a theorem to be delegated to a
library theory solver.

\subsection{Qualifier family: sorting semantics}
The system introduces qualifiers corresponding to \(\mathsf{sorted}(ys)\) and \(\mathsf{perm}(ys,xs)\) for order-sensitive library routines. In a standard
Liquid-style design, such qualifiers would be mined from syntax and then used
to instantiate implication constraints. In our expanded setting, they are also
fed into the global optimization layer, where they compete with richer
dependent artifacts such as protocol summaries, witness-returning signatures,
and library-theory lemmas. This is especially important in Python, because a
user may write no annotations at all, yet the code still contains local uses
that force semantic obligations. The system should therefore create the
candidate even when it has not yet decided whether to surface it as a user
precondition, a hidden proof obligation, or a theorem to be delegated to a
library theory solver.

\subsection{Qualifier family: numeric safety}
The system introduces qualifiers corresponding to \(d \neq 0\) for division and modulo, \(n \ge 0\) for natural indices and lengths. In a standard
Liquid-style design, such qualifiers would be mined from syntax and then used
to instantiate implication constraints. In our expanded setting, they are also
fed into the global optimization layer, where they compete with richer
dependent artifacts such as protocol summaries, witness-returning signatures,
and library-theory lemmas. This is especially important in Python, because a
user may write no annotations at all, yet the code still contains local uses
that force semantic obligations. The system should therefore create the
candidate even when it has not yet decided whether to surface it as a user
precondition, a hidden proof obligation, or a theorem to be delegated to a
library theory solver.



\section{Unannotated programs as optimization instances}

The zero-annotation case is not a corner case; it is the default case we should
expect in ordinary Python repositories. That means every unannotated function
must still induce a meaningful optimization problem.

Consider the completely unannotated program:
\begin{lstlisting}[style=deppy,language=Python]
def first(xs):
    return xs[0]
\end{lstlisting}
A purely annotation-driven system learns almost nothing here. By contrast, an
optimization-based dependent system can still derive useful obligations.

\begin{enumerate}[label=\arabic*.]
  \item The indexing expression generates a safety obligation suggesting
  \(\mathrm{len}(xs) > 0\).
  \item The variable \py{xs} becomes a candidate for a non-empty collection
  contract.
  \item The return position may inherit an existentially unknown element type,
  yielding a latent family \(\exists \alpha.\; \mathsf{List}(\alpha) \to \alpha\)
  plus the precondition \(\mathrm{len}(xs) > 0\).
  \item The optimizer decides whether to surface this as a minimal
  \py{@deppy.requires("len(xs) > 0")} decorator suggestion, a hidden internal
  precondition, or a proof obligation attached to the function summary.
\end{enumerate}

This is a useful point of departure from traditional type inference. The system
is not merely searching for a principal Hindley--Milner type. It is searching
for the best explanatory dependent contract it can justify from the program
surface plus the theory inventory available to it.

\begin{proposition}[Obligation synthesis from dangerous operations]
Let $e$ be an expression that is only safe when some predicate $\phi$ holds,
such as indexing, division, attribute access on an optional value, or calling a
protocol method. Then the typing engine may soundly generate $\phi$ as a
candidate precondition for the containing function, provided the obligation is
scoped to the relevant inputs and the generated contract is clearly classified as
proved, assumed, bounded, or unknown.
\end{proposition}



\section{Formal optimization model}

We now present a more explicit mathematical formulation. Let $P$ be the set of
program points in a module, $V$ the set of source variables and heap locations,
and $C$ the set of candidate dependent information objects. We partition $C$ as
\[
C = C_{\mathrm{ref}} \uplus C_{\Pi} \uplus C_{\Sigma} \uplus C_{\mathrm{eq}} \uplus C_{\mathrm{heap}} \uplus C_{\mathrm{theory}} \uplus C_{\mathrm{proto}}.
\]
The components correspond respectively to refinements, dependent functions,
dependent pairs, equality claims, heap summaries, library-theory facts, and
protocol/structural interface facts.

Each candidate $c$ carries:
\[
\mathsf{meta}(c)=\langle
\support(c),
\trust(c),
\cost(c),
\mathsf{frag}(c),
\mathsf{coverage}(c),
\mathsf{burden}(c),
\mathsf{origin}(c)
\rangle.
\]
A selection is a Boolean vector $z \in \{0,1\}^{|C|}$. We introduce auxiliary
coverage variables $u_v$ for variables or semantic slots covered by the selected
facts. Then a representative objective is:
\[
\max_{z,u}
\left(
\sum_{v \in V} \alpha_v u_v
+
\sum_{c \in C} z_c\big(
\beta_1 \support(c)
+
\beta_2 \mathsf{expressivity}(c)
+
\beta_3 \mathsf{reuse}(c)
-
\beta_4 \cost(c)
-
\beta_5 \mathsf{burden}(c)
\big)
-
\sum_{(c,c') \in E_{\mathrm{overlap}}} \gamma_{cc'} z_c z_{c'}
\right)
\]
subject to:
\begin{align*}
&u_v \Leftrightarrow \bigvee_{c : v \in \mathsf{coverage}(c)} z_c, \\
&\sum_{c \in C} z_c\,\mathsf{burden}(c) \le B, \\
&\sum_{c \in C} z_c\,\mathsf{solverWeight}(c) \le S, \\
&z_c = 0 \text{ whenever } \mathsf{frag}(c) \text{ is excluded by the current trusted mode.}
\end{align*}

This formulation explains why \Zthree{}-backed optimization is natural. The
problem is not simply a yes/no satisfiability check; it is a weighted selection
problem with hard and soft constraints. MaxSMT, Optimize, or related 0--1
integer formulations are therefore a direct fit for practical implementations.



\section{A CEGIS-class framework for dependent typing: CEGDTS}

We now introduce a named synthesis framework intended to play, for Pythonic
dependent typing, a role analogous to the role that CEGIS plays for program
synthesis. We call it \emph{CounterExample-Guided Dependent Type Synthesis}
(\textsc{CEGDTS}). The basic insight is that dependent typing in a low-
annotation, theory-rich language is not only a checking problem and not only an
optimization problem. It is a \emph{synthesis} problem over typing artifacts.
The system must synthesize a contract, qualifier set, witness structure,
protocol summary, heap summary, and theory instantiation that is strong enough
to justify the program's intended uses, weak enough to remain automatable, and
compact enough to be presented back to a user or LLM in a useful form.

Classical CEGIS alternates between synthesis and verification:
\[
\text{candidate} \longrightarrow \text{verify} \longrightarrow \text{counterexample} \longrightarrow \text{refine synthesis space}.
\]
For dependent typing, we require a richer state because the synthesizer is not
searching over whole programs but over semantic typing objects. We therefore
write the synthesis state as
\[
\Sigma_i = \langle C_i, Q_i, T_i, W_i, H_i, \Theta_i \rangle
\]
where:
\begin{itemize}
  \item $C_i$ is the candidate contract family,
  \item $Q_i$ is the qualifier bank,
  \item $T_i$ is the activated library-theory set,
  \item $W_i$ is the witness schema family (for Sigma types, equality objects,
  and theory certificates),
  \item $H_i$ is the heap-and-protocol summary family,
  \item $\Theta_i$ is the current optimization weight vector or policy state.
\end{itemize}

The synthesizer proposes a typing package
\[
S_i \in \mathcal{S}(\Sigma_i)
\]
by solving an optimization-and-synthesis problem:
\[
S_i = \arg\max_{S \in \mathcal{S}(\Sigma_i)}
\Big(
\mathsf{Info}(S)
- \lambda_1 \mathsf{SolveCost}(S)
- \lambda_2 \mathsf{AnnotationCost}(S)
- \lambda_3 \mathsf{TheoryRisk}(S)
- \lambda_4 \mathsf{HeapInstability}(S)
\Big)
\]
subject to hard typing compatibility constraints. The verifier then checks
\[
\mathsf{Verify}(M, S_i) \in \{\Proved, \Refuted, \Unknown, \Assumed, \Bounded\}
\]
and, more importantly, returns structured evidence:
\[
E_i = \langle \kappa_i, \pi_i, u_i, \rho_i \rangle
\]
where $\kappa_i$ is a counterexample model when available, $\pi_i$ is a proof
object or unsat-core explanation when available, $u_i$ is usage telemetry such
as which obligations consumed the most solver budget, and $\rho_i$ is a repair
profile describing which missing facts would have blocked the failure.

The central novelty is that the refinement step is no longer merely a predicate
addition function. It is a typed synthesis-space transformer:
\[
\Sigma_{i+1} = \mathcal{R}(\Sigma_i, E_i)
\]
with components:
\begin{align*}
C_{i+1} &= C_i \cup \mathsf{ExpandContracts}(E_i), \\
Q_{i+1} &= Q_i \cup \mathsf{ExpandQualifiers}(E_i), \\
T_{i+1} &= T_i \cup \mathsf{ActivateTheories}(E_i), \\
W_{i+1} &= W_i \cup \mathsf{ExpandWitnessSchemas}(E_i), \\
H_{i+1} &= H_i \cup \mathsf{RefineHeapSummaries}(E_i), \\
\Theta_{i+1} &= \mathsf{UpdateWeights}(\Theta_i, E_i).
\end{align*}
This refinement operator is what makes \textsc{CEGDTS} genuinely richer than a
plain qualifier loop. A failure at a PyTorch sort example, for instance, need
not merely add a scalar predicate. It may activate a permutation theory,
introduce a witness schema for sortedness certificates, or lower the score of an
annotation-expensive contract family in favor of a more reusable theory-backed
summary.

\begin{definition}[\textsc{CEGDTS} step]
A single \textsc{CEGDTS} step is the transition
\[
\Sigma_i \xrightarrow{\ \mathsf{Synthesize}\ }
S_i \xrightarrow{\ \mathsf{Verify}\ }
E_i \xrightarrow{\ \mathsf{Repair}\ }
\Sigma_{i+1}.
\]
The framework is \emph{soundly staged} when every synthesized package $S_i$ is
verified only through the explicit trusted kernel, and every unproved component
of $S_i$ is preserved as a visible proof obligation.
\end{definition}

\begin{definition}[Repair completeness up to a candidate basis]
Fix a finite basis of possible contracts, qualifiers, theory activations,
witness schemas, and heap summaries. A repair operator $\mathcal{R}$ is
\emph{complete up to the basis} if, whenever some basis element would eliminate a
current counterexample or prove a currently unknown obligation, repeated
application of $\mathcal{R}$ eventually introduces that basis element into the
search state.
\end{definition}

\begin{theorem}[Finite-basis stabilization of \textsc{CEGDTS}]
Assume a finite candidate basis, a sound verifier, and a repair operator that is
complete up to the basis. Then the \textsc{CEGDTS} loop eventually stabilizes in
one of three states:
\begin{enumerate}[label=\arabic*.]
  \item a package whose trusted obligations are proved,
  \item a refuted package together with a stable counterexample pattern,
  \item a frontier package in which every remaining improvement would require
  leaving the current basis or trusted fragment.
\end{enumerate}
\end{theorem}

\begin{proof}[Proof sketch]
Because the basis is finite, there are only finitely many synthesis states up to
weight updates and selection equivalence. Completeness up to the basis implies
that any eliminable failure mode is eventually repaired. Sound verification
prevents spurious convergence to false proofs. Therefore the loop cannot refine
strictly forever without exhausting the finite basis; it must stabilize at a
proved, refuted, or frontier state.
\end{proof}

\begin{remark}
The frontier case is the crucial one for a practical dependent system. It is the
formal place where the implementation can say: ``we have inferred the strongest
contract currently supported by our candidate basis, trusted fragment, and
library theories; here are the exact residual proof goals.'' In this sense,
\textsc{CEGDTS} is not just a synthesis algorithm but a discipline for exposing
what remains semantically missing.
\end{remark}

\subsection{Why this matters for normal Python}
The synthesis view is especially natural for ordinary Pythonic code, because the
user often writes only part of the final specification. A class method might
implicitly require a protocol; an index operation might imply a length bound; a
PyTorch call might require a shape/device theorem; a sorting routine might need
permutation and monotonicity obligations. \textsc{CEGDTS} explains how such facts
can be synthesized, verified, and repaired iteratively without forcing the user
to write them all up front.

\subsection{Why this is not merely refinement inference}
Refinement inference alone would update $Q_i$. \textsc{CEGDTS} updates the full
state $\Sigma_i$, including witness schemas $W_i$, heap summaries $H_i$, and
library theories $T_i$. This is why it is a candidate for a CEGIS-class idea for
Pythonic dependent typing rather than just another qualifier-generation trick.


\section{Library support as new theories}

One of the most important practical departures from classical refinement typing
is that library support should be modeled as an extensible family of semantic
theories rather than as ad hoc handwritten exceptions. The system should treat
each serious library domain as a \emph{theory module} with its own symbols,
obligation generators, trusted kernels, and optimization costs.
\subsection{Theory interface for PyTorch}
The PyTorch theory should provide relations for tensor rank, shape equalities, broadcasting, device agreement, batch discipline, sorting and top-k semantics, permutation and monotonicity obligations, view/reshape side conditions. From the perspective of
the core inference engine, the theory supplies three things: candidate facts,
proof procedures or reductions into the trusted core, and an explicit cost model.
This design permits the optimizer to reason not only about generic qualifiers but
also about whether adding a library-theoretic fact is worth its solver and trust
cost. In particular, PyTorch is not just ``a module import. It contributes a
shape/device/permutation theory that may generate dependent contracts over tensor
programs while still preserving normal Python source syntax.

\subsection{Theory interface for NumPy}
The NumPy theory should provide relations for shape algebra, broadcasting, dtype stability, axis semantics, memory-layout side conditions, reduction output ranks. From the perspective of
the core inference engine, the theory supplies three things: candidate facts,
proof procedures or reductions into the trusted core, and an explicit cost model.
This design permits the optimizer to reason not only about generic qualifiers but
also about whether adding a library-theoretic fact is worth its solver and trust
cost. In particular, PyTorch is not just ``a module import. It contributes a
shape/device/permutation theory that may generate dependent contracts over tensor
programs while still preserving normal Python source syntax.

\subsection{Theory interface for Pandas}
The Pandas theory should provide relations for column presence, schema preservation, index alignment, grouping keys, merge compatibility, missing-value propagation. From the perspective of
the core inference engine, the theory supplies three things: candidate facts,
proof procedures or reductions into the trusted core, and an explicit cost model.
This design permits the optimizer to reason not only about generic qualifiers but
also about whether adding a library-theoretic fact is worth its solver and trust
cost. In particular, PyTorch is not just ``a module import. It contributes a
shape/device/permutation theory that may generate dependent contracts over tensor
programs while still preserving normal Python source syntax.

\subsection{Theory interface for SQLAlchemy}
The SQLAlchemy theory should provide relations for table schemas, field nullability, key relations, query result arity, transaction effects. From the perspective of
the core inference engine, the theory supplies three things: candidate facts,
proof procedures or reductions into the trusted core, and an explicit cost model.
This design permits the optimizer to reason not only about generic qualifiers but
also about whether adding a library-theoretic fact is worth its solver and trust
cost. In particular, PyTorch is not just ``a module import. It contributes a
shape/device/permutation theory that may generate dependent contracts over tensor
programs while still preserving normal Python source syntax.

\subsection{Theory interface for asyncio}
The asyncio theory should provide relations for awaitable protocol conformance, task lifecycle states, queue non-emptiness, cancellation obligations. From the perspective of
the core inference engine, the theory supplies three things: candidate facts,
proof procedures or reductions into the trusted core, and an explicit cost model.
This design permits the optimizer to reason not only about generic qualifiers but
also about whether adding a library-theoretic fact is worth its solver and trust
cost. In particular, PyTorch is not just ``a module import. It contributes a
shape/device/permutation theory that may generate dependent contracts over tensor
programs while still preserving normal Python source syntax.

\subsection{Theory interface for dataclasses and pydantic}
The dataclasses and pydantic theory should provide relations for field initialization, optionality, validation postconditions, constructor-produced invariants. From the perspective of
the core inference engine, the theory supplies three things: candidate facts,
proof procedures or reductions into the trusted core, and an explicit cost model.
This design permits the optimizer to reason not only about generic qualifiers but
also about whether adding a library-theoretic fact is worth its solver and trust
cost. In particular, PyTorch is not just ``a module import. It contributes a
shape/device/permutation theory that may generate dependent contracts over tensor
programs while still preserving normal Python source syntax.

\subsection{Theory interface for collections and iterators}
The collections and iterators theory should provide relations for non-emptiness, iterator exhaustion, length-sensitive indexing, partitioning and chunking invariants. From the perspective of
the core inference engine, the theory supplies three things: candidate facts,
proof procedures or reductions into the trusted core, and an explicit cost model.
This design permits the optimizer to reason not only about generic qualifiers but
also about whether adding a library-theoretic fact is worth its solver and trust
cost. In particular, PyTorch is not just ``a module import. It contributes a
shape/device/permutation theory that may generate dependent contracts over tensor
programs while still preserving normal Python source syntax.



\section{Classes, protocols, heap modeling, and normal object-oriented Python}

Supporting normal Python requires native treatment of classes, protocols, and
heap updates. A function-only calculus is insufficient. We therefore model:
\begin{itemize}
  \item nominal class types with declared and inferred fields,
  \item protocol types representing structural interfaces,
  \item heap write summaries generated from assignments such as \py{self.x = e},
  \item method contracts that may refer to object state,
  \item structural obligations induced by attribute access and method calls,
  \item heap-sensitive invariants that can be surfaced as dependent contracts.
\end{itemize}

A class declaration contributes both nominal and structural information. The
nominal part matters for constructor-typed call checking and inheritance; the
structural part matters for protocol satisfaction and field reasoning. This dual
view is especially important in Python because ordinary code frequently mixes
nominal idioms with duck typing. The heap summary then records where state is
written and which field types are expected to persist.

\begin{example}[Heap-sensitive invariant]
Consider
\begin{lstlisting}[style=deppy,language=Python]
class Counter:
    value: int

    def __init__(self) -> None:
        self.value = 0
\end{lstlisting}
The field declaration and constructor write together induce a persistent heap
invariant that \py{self.value} should remain an \py{int}. A later write
\py{self.value = "zero"} should therefore be flagged as a static mismatch, even
if the surrounding source is otherwise unannotated.
\end{example}



\section{Ordinary Python examples the system should catch, accept, or defer}

The example corpus is not an afterthought. It is the operational meaning of the
design. We therefore distinguish three categories.

\subsection{Category A: bugs that current static analysis should catch}
Examples include:
\begin{itemize}
  \item returning a string from a function whose inferred or declared result is
  numeric or vector-valued,
  \item violating a protocol requirement by passing a class whose method returns
  the wrong type,
  \item writing a value of the wrong field type to a known heap location,
  \item calling a function with an argument whose static type does not satisfy the
  expected nominal or protocol contract,
  \item obviously violating a dependent signature by returning a completely
  incompatible sort of object.
\end{itemize}

\subsection{Category B: scripts that should be accepted with few or no annotations}
Examples include:
\begin{itemize}
  \item guard-based head functions where the system can infer a non-emptiness
  precondition,
  \item class wrappers that simply preserve declared field types,
  \item functions whose return type is a simple arithmetic or structural fact that
  can be inferred from operations such as \py{len}, literal construction, or
  constructor calls,
  \item theory-backed library uses where the trusted core can prove the relevant
  safety obligations.
\end{itemize}

\subsection{Category C: scripts whose meaning is expressible but not fully auto-proved}
This is where the system must remain honest. Sorting correctness, stability of
custom comparators, exact tensor permutation semantics under exotic library
transformations, or semantic data-structure invariants may be expressible as
proof goals even when the current automatic core cannot discharge them. Such
examples must not be discarded. They are part of what it means to support full
dependent types instead of only a thin verified fragment.



\section{PyTorch as a flagship theory: shapes, devices, and sorting correctness}

PyTorch is a particularly important test case because it demonstrates why
library theories matter. A typical PyTorch user writes code that looks like
normal Python classes and methods, often over \py{nn.Module}, with almost no
extra specification syntax. Yet the semantic content is highly structured: shape
relations, device consistency, rank transformations, monotonicity properties,
index/value alignment, and in some cases full algorithmic correctness.

A TensorGuard-style model already captures some of this landscape at the level
of shapes and devices. The broader dependent design proposed here pushes farther.
It treats library semantics as a theory capable of supplying richer propositions
such as:
\[
\mathsf{sorted}(y), \qquad \mathsf{perm}(y,x), \qquad
\mathsf{shape}(y)=\mathsf{shape}(x), \qquad
\mathsf{device}(y)=\mathsf{device}(x).
\]

\subsection{A TensorGuard-style PyTorch sorting module}
A deliberately Pythonic example is:
\begin{lstlisting}[style=deppy,language=Python]
import torch
import torch.nn as nn

class SortingModule(nn.Module):
    @deppy.guarantee({
        "parameters": [{"name": "xs", "type": "Vec(Int, n)"}],
        "return": "Vec(Int, n)",
        "signature": "(n : Nat) -> (xs : Vec(Int, n)) -> Vec(Int, n)",
        "proof_goals": ["sorted(result)", "perm(result, xs)"]
    })
    def forward(self, xs):
        values, _indices = torch.sort(xs, dim=-1)
        return values
\end{lstlisting}
This example is important for three reasons. First, it looks like ordinary
PyTorch code rather than a theorem prover term. Second, it carries a full
dependent contract, not only a scalar refinement. Third, it illustrates the
honest trust boundary: the result shape preservation may be dischargeable by the
trusted core, while semantic goals such as sortedness and permutation may remain
as explicit proof obligations if the current theory module does not yet contain
all required lemmas.

\subsection{A static bug in a PyTorch sorting module}
Now consider the buggy variant:
\begin{lstlisting}[style=deppy,language=Python]
class BuggySortingModule(nn.Module):
    @deppy.guarantee({
        "parameters": [{"name": "xs", "type": "Vec(Int, n)"}],
        "return": "Vec(Int, n)",
        "signature": "(n : Nat) -> (xs : Vec(Int, n)) -> Vec(Int, n)",
        "proof_goals": ["sorted(result)", "perm(result, xs)"]
    })
    def forward(self, xs):
        values, _indices = torch.sort(xs, dim=-1)
        _ = values
        return "not-a-vector"
\end{lstlisting}
Here the dependent specification is semantically rich, yet the static bug is
simple: the returned object is the wrong kind entirely. This is exactly the sort
of example a practical system should catch immediately. Notice the interplay:
full dependent expressivity is present, but the static analysis can still deliver
a low-level type mismatch without having to solve the deepest semantic proof
goals first.



\section{Algorithms for extraction, scoring, and contract compression}

A workable implementation can be described as the following staged algorithm.

\begin{enumerate}[label=\arabic*.]
  \item Parse ordinary Python into a native-Python intermediate form that retains
  classes, protocols, heap writes, control-flow structure, and relevant library
  calls.
  \item Harvest guards, assertions, protocol uses, attribute requirements, index
  safety obligations, and theory-generated facts.
  \item Generate qualifier candidates in the \LiquidHaskell{} style, but broaden
  the candidate family to include dependent signatures, witness-returning
  contracts, protocol summaries, heap invariants, and theory lemmas.
  \item Score or optimize over these candidates using a solver-backed objective.
  \item Lower the selected fragment to a trusted core for \Zthree{} or other
  decision procedures.
  \item Preserve unlowered obligations as explicit proof goals.
  \item Compress the chosen information into user-facing outputs: inferred
  contracts, minimal annotation suggestions, diagnostics, and proof artifacts.
\end{enumerate}

Although the implementation is staged rather than monolithic, the logic can be
summarized abstractly as:
\[
\mathsf{Infer}(M) = \mathsf{Compress}\big(
  \mathsf{Select}(\mathsf{Candidates}(\mathsf{Extract}(M)))
\big).
\]
The crucial point is that \(\mathsf{Select}\) is not a trivial filter. It is the
optimization heart of the system.



\section{Trusted kernel, partial automation, and explicit proof obligations}

The trusted boundary should be visible in the result type of the checker. Every
obligation should end in one of a small number of explicit verdict states,
including \Proved{}, \Refuted{}, \Bounded{}, \Assumed{}, and \Unknown{}. This is
how the design reconciles ambitious dependent expressivity with practical
solver-backed automation.

\begin{definition}[Verdict semantics]
Given an obligation $o$, the checker returns:
\begin{itemize}
  \item \Proved{} if $o$ is discharged within the trusted kernel,
  \item \Refuted{} if a trusted counterexample is found,
  \item \Bounded{} if $o$ holds under an explicitly finite approximation,
  \item \Assumed{} if $o$ is retained under an unproved but visible assumption,
  \item \Unknown{} if the current procedures neither prove nor refute $o$.
\end{itemize}
\end{definition}

This classification is especially important for full dependent types over
Pythonic source. Once we allow theory-indexed obligations such as sorting
correctness, heap ownership, or precise tensor permutation semantics, it is no
longer acceptable to force a false binary choice between ``verified'' and
``untyped.'' The verdict algebra is part of the language design.


\section{Extended case-study catalogue}

To keep the document operational rather than merely aspirational, we close with
a catalogue of case-study shapes that the implementation should support.
\subsection{Protocol mismatch}
A class method returning the wrong type for a protocol requirement should yield a direct call-argument mismatch.
The purpose of this example family is twofold. First, it anchors the design in
code that still reads like Python. Second, it creates an evaluation harness in
which the system must decide whether the example should be accepted, rejected,
or carried forward with explicit proof obligations. In an optimization-first
design, each example also induces a different budget tradeoff: some prefer
simple local qualifiers, others require theory lemmas, and others demand that
the checker preserve semantic obligations without pretending to have proved
them. This is precisely the behavior we want from a practical dependent system
for Pythonic code.

\subsection{Heap field mismatch}
Writing a string into a field declared as integer should violate the persistent heap invariant inferred from the class declaration.
The purpose of this example family is twofold. First, it anchors the design in
code that still reads like Python. Second, it creates an evaluation harness in
which the system must decide whether the example should be accepted, rejected,
or carried forward with explicit proof obligations. In an optimization-first
design, each example also induces a different budget tradeoff: some prefer
simple local qualifiers, others require theory lemmas, and others demand that
the checker preserve semantic obligations without pretending to have proved
them. This is precisely the behavior we want from a practical dependent system
for Pythonic code.

\subsection{Guarded head function}
A low-annotation function using a non-empty guard should be accepted and compressed into a minimal precondition.
The purpose of this example family is twofold. First, it anchors the design in
code that still reads like Python. Second, it creates an evaluation harness in
which the system must decide whether the example should be accepted, rejected,
or carried forward with explicit proof obligations. In an optimization-first
design, each example also induces a different budget tradeoff: some prefer
simple local qualifiers, others require theory lemmas, and others demand that
the checker preserve semantic obligations without pretending to have proved
them. This is precisely the behavior we want from a practical dependent system
for Pythonic code.

\subsection{Unannotated indexing function}
A fully unannotated function returning xs[0] should synthesize a non-emptiness obligation rather than remaining opaque.
The purpose of this example family is twofold. First, it anchors the design in
code that still reads like Python. Second, it creates an evaluation harness in
which the system must decide whether the example should be accepted, rejected,
or carried forward with explicit proof obligations. In an optimization-first
design, each example also induces a different budget tradeoff: some prefer
simple local qualifiers, others require theory lemmas, and others demand that
the checker preserve semantic obligations without pretending to have proved
them. This is precisely the behavior we want from a practical dependent system
for Pythonic code.

\subsection{Constructor misuse}
Passing a string where a class-typed parameter is required should be rejected even when the constructor itself has no dependent annotation.
The purpose of this example family is twofold. First, it anchors the design in
code that still reads like Python. Second, it creates an evaluation harness in
which the system must decide whether the example should be accepted, rejected,
or carried forward with explicit proof obligations. In an optimization-first
design, each example also induces a different budget tradeoff: some prefer
simple local qualifiers, others require theory lemmas, and others demand that
the checker preserve semantic obligations without pretending to have proved
them. This is precisely the behavior we want from a practical dependent system
for Pythonic code.

\subsection{PyTorch shape bug}
A TensorGuard-style model with a wrong linear dimension should remain representable as a theory-backed example, even if this surface layer does not yet prove the shape theorem automatically.
The purpose of this example family is twofold. First, it anchors the design in
code that still reads like Python. Second, it creates an evaluation harness in
which the system must decide whether the example should be accepted, rejected,
or carried forward with explicit proof obligations. In an optimization-first
design, each example also induces a different budget tradeoff: some prefer
simple local qualifiers, others require theory lemmas, and others demand that
the checker preserve semantic obligations without pretending to have proved
them. This is precisely the behavior we want from a practical dependent system
for Pythonic code.

\subsection{PyTorch sorting correctness}
A sorting module should preserve a shape/index family and carry proof goals for sortedness and permutation.
The purpose of this example family is twofold. First, it anchors the design in
code that still reads like Python. Second, it creates an evaluation harness in
which the system must decide whether the example should be accepted, rejected,
or carried forward with explicit proof obligations. In an optimization-first
design, each example also induces a different budget tradeoff: some prefer
simple local qualifiers, others require theory lemmas, and others demand that
the checker preserve semantic obligations without pretending to have proved
them. This is precisely the behavior we want from a practical dependent system
for Pythonic code.

\subsection{Sigma witness export}
A witness-returning function should demonstrate that the implementation supports dependent pairs and not only refinements.
The purpose of this example family is twofold. First, it anchors the design in
code that still reads like Python. Second, it creates an evaluation harness in
which the system must decide whether the example should be accepted, rejected,
or carried forward with explicit proof obligations. In an optimization-first
design, each example also induces a different budget tradeoff: some prefer
simple local qualifiers, others require theory lemmas, and others demand that
the checker preserve semantic obligations without pretending to have proved
them. This is precisely the behavior we want from a practical dependent system
for Pythonic code.


\appendix
\section{Design checklist}

For completeness, we record a checklist of the principles that recur throughout this file.
\begin{enumerate}[label=\arabic*.]
  \item Keep code visually close to normal Python.
  \item Treat annotations as optional evidence rather than mandatory syntax.
  \item Support full dependent structure, not only refinement predicates.
  \item Model classes, protocols, heaps, and library semantics directly.
  \item Treat inference as optimization, not mere derivability.
  \item Use \Zthree{} and MaxSMT/Optimize for the decidable trusted core.
  \item Retain rich obligations as visible proof goals when automation stops.
  \item Exploit Liquid-style qualifier mining, but enlarge it to a broader candidate language.
  \item Use library theories, especially PyTorch, as first-class semantic modules.
  \item Prefer minimal user-facing annotations even when the internal proof state is rich.
  \item Keep code visually close to normal Python.
  \item Treat annotations as optional evidence rather than mandatory syntax.
  \item Support full dependent structure, not only refinement predicates.
  \item Model classes, protocols, heaps, and library semantics directly.
  \item Treat inference as optimization, not mere derivability.
  \item Use \Zthree{} and MaxSMT/Optimize for the decidable trusted core.
  \item Retain rich obligations as visible proof goals when automation stops.
  \item Exploit Liquid-style qualifier mining, but enlarge it to a broader candidate language.
  \item Use library theories, especially PyTorch, as first-class semantic modules.
  \item Prefer minimal user-facing annotations even when the internal proof state is rich.
  \item Keep code visually close to normal Python.
  \item Treat annotations as optional evidence rather than mandatory syntax.
  \item Support full dependent structure, not only refinement predicates.
  \item Model classes, protocols, heaps, and library semantics directly.
  \item Treat inference as optimization, not mere derivability.
  \item Use \Zthree{} and MaxSMT/Optimize for the decidable trusted core.
  \item Retain rich obligations as visible proof goals when automation stops.
  \item Exploit Liquid-style qualifier mining, but enlarge it to a broader candidate language.
  \item Use library theories, especially PyTorch, as first-class semantic modules.
  \item Prefer minimal user-facing annotations even when the internal proof state is rich.
  \item Keep code visually close to normal Python.
  \item Treat annotations as optional evidence rather than mandatory syntax.
  \item Support full dependent structure, not only refinement predicates.
  \item Model classes, protocols, heaps, and library semantics directly.
  \item Treat inference as optimization, not mere derivability.
  \item Use \Zthree{} and MaxSMT/Optimize for the decidable trusted core.
  \item Retain rich obligations as visible proof goals when automation stops.
  \item Exploit Liquid-style qualifier mining, but enlarge it to a broader candidate language.
  \item Use library theories, especially PyTorch, as first-class semantic modules.
  \item Prefer minimal user-facing annotations even when the internal proof state is rich.
  \item Keep code visually close to normal Python.
  \item Treat annotations as optional evidence rather than mandatory syntax.
  \item Support full dependent structure, not only refinement predicates.
  \item Model classes, protocols, heaps, and library semantics directly.
  \item Treat inference as optimization, not mere derivability.
  \item Use \Zthree{} and MaxSMT/Optimize for the decidable trusted core.
  \item Retain rich obligations as visible proof goals when automation stops.
  \item Exploit Liquid-style qualifier mining, but enlarge it to a broader candidate language.
  \item Use library theories, especially PyTorch, as first-class semantic modules.
  \item Prefer minimal user-facing annotations even when the internal proof state is rich.
  \item Keep code visually close to normal Python.
  \item Treat annotations as optional evidence rather than mandatory syntax.
  \item Support full dependent structure, not only refinement predicates.
  \item Model classes, protocols, heaps, and library semantics directly.
  \item Treat inference as optimization, not mere derivability.
  \item Use \Zthree{} and MaxSMT/Optimize for the decidable trusted core.
  \item Retain rich obligations as visible proof goals when automation stops.
  \item Exploit Liquid-style qualifier mining, but enlarge it to a broader candidate language.
  \item Use library theories, especially PyTorch, as first-class semantic modules.
  \item Prefer minimal user-facing annotations even when the internal proof state is rich.
  \item Keep code visually close to normal Python.
  \item Treat annotations as optional evidence rather than mandatory syntax.
  \item Support full dependent structure, not only refinement predicates.
  \item Model classes, protocols, heaps, and library semantics directly.
  \item Treat inference as optimization, not mere derivability.
  \item Use \Zthree{} and MaxSMT/Optimize for the decidable trusted core.
  \item Retain rich obligations as visible proof goals when automation stops.
  \item Exploit Liquid-style qualifier mining, but enlarge it to a broader candidate language.
  \item Use library theories, especially PyTorch, as first-class semantic modules.
  \item Prefer minimal user-facing annotations even when the internal proof state is rich.
  \item Keep code visually close to normal Python.
  \item Treat annotations as optional evidence rather than mandatory syntax.
  \item Support full dependent structure, not only refinement predicates.
  \item Model classes, protocols, heaps, and library semantics directly.
  \item Treat inference as optimization, not mere derivability.
  \item Use \Zthree{} and MaxSMT/Optimize for the decidable trusted core.
  \item Retain rich obligations as visible proof goals when automation stops.
  \item Exploit Liquid-style qualifier mining, but enlarge it to a broader candidate language.
  \item Use library theories, especially PyTorch, as first-class semantic modules.
  \item Prefer minimal user-facing annotations even when the internal proof state is rich.
\end{enumerate}

\section{A final statement of emphasis}
The most important emphasis of this monograph can be stated in one sentence:
\DepPy{} should let programmers and LLMs write ordinary Python-looking code,
sometimes with no annotations at all, while recovering the strongest practical
dependent contracts and proof obligations that the current solver-theoretic and
library-theoretic budget can justify. The optimization formulation is not a
side feature; it is the organizing principle that lets ordinary syntax, sparse
annotations, full dependent expressivity, theory-specific reasoning, and an
explicit trust boundary coexist in one system.
\section{Supplementary technical elaborations}
\subsection{Supplement: constraint normalization}
This supplementary note elaborates the role of constraint normalization in the overall
optimization-based dependent typing architecture. The recurring theme is that
Pythonic source programs expose many plausible semantic candidates, but the
system must decide which of them should become trusted core obligations, which
should be surfaced as suggested annotations, which should remain as proof
goals, and which should be discarded as too unstable, too expensive, or too
poorly supported. This makes the architecture simultaneously a typing system,
a semantic extraction engine, a theory-integration framework, and an explicit
optimization procedure over candidate dependent information objects.

This supplementary note elaborates the role of constraint normalization in the overall
optimization-based dependent typing architecture. The recurring theme is that
Pythonic source programs expose many plausible semantic candidates, but the
system must decide which of them should become trusted core obligations, which
should be surfaced as suggested annotations, which should remain as proof
goals, and which should be discarded as too unstable, too expensive, or too
poorly supported. This makes the architecture simultaneously a typing system,
a semantic extraction engine, a theory-integration framework, and an explicit
optimization procedure over candidate dependent information objects.

This supplementary note elaborates the role of constraint normalization in the overall
optimization-based dependent typing architecture. The recurring theme is that
Pythonic source programs expose many plausible semantic candidates, but the
system must decide which of them should become trusted core obligations, which
should be surfaced as suggested annotations, which should remain as proof
goals, and which should be discarded as too unstable, too expensive, or too
poorly supported. This makes the architecture simultaneously a typing system,
a semantic extraction engine, a theory-integration framework, and an explicit
optimization procedure over candidate dependent information objects.

This supplementary note elaborates the role of constraint normalization in the overall
optimization-based dependent typing architecture. The recurring theme is that
Pythonic source programs expose many plausible semantic candidates, but the
system must decide which of them should become trusted core obligations, which
should be surfaced as suggested annotations, which should remain as proof
goals, and which should be discarded as too unstable, too expensive, or too
poorly supported. This makes the architecture simultaneously a typing system,
a semantic extraction engine, a theory-integration framework, and an explicit
optimization procedure over candidate dependent information objects.

This supplementary note elaborates the role of constraint normalization in the overall
optimization-based dependent typing architecture. The recurring theme is that
Pythonic source programs expose many plausible semantic candidates, but the
system must decide which of them should become trusted core obligations, which
should be surfaced as suggested annotations, which should remain as proof
goals, and which should be discarded as too unstable, too expensive, or too
poorly supported. This makes the architecture simultaneously a typing system,
a semantic extraction engine, a theory-integration framework, and an explicit
optimization procedure over candidate dependent information objects.

This supplementary note elaborates the role of constraint normalization in the overall
optimization-based dependent typing architecture. The recurring theme is that
Pythonic source programs expose many plausible semantic candidates, but the
system must decide which of them should become trusted core obligations, which
should be surfaced as suggested annotations, which should remain as proof
goals, and which should be discarded as too unstable, too expensive, or too
poorly supported. This makes the architecture simultaneously a typing system,
a semantic extraction engine, a theory-integration framework, and an explicit
optimization procedure over candidate dependent information objects.

\subsection{Supplement: proof goal retention}
This supplementary note elaborates the role of proof goal retention in the overall
optimization-based dependent typing architecture. The recurring theme is that
Pythonic source programs expose many plausible semantic candidates, but the
system must decide which of them should become trusted core obligations, which
should be surfaced as suggested annotations, which should remain as proof
goals, and which should be discarded as too unstable, too expensive, or too
poorly supported. This makes the architecture simultaneously a typing system,
a semantic extraction engine, a theory-integration framework, and an explicit
optimization procedure over candidate dependent information objects.

This supplementary note elaborates the role of proof goal retention in the overall
optimization-based dependent typing architecture. The recurring theme is that
Pythonic source programs expose many plausible semantic candidates, but the
system must decide which of them should become trusted core obligations, which
should be surfaced as suggested annotations, which should remain as proof
goals, and which should be discarded as too unstable, too expensive, or too
poorly supported. This makes the architecture simultaneously a typing system,
a semantic extraction engine, a theory-integration framework, and an explicit
optimization procedure over candidate dependent information objects.

This supplementary note elaborates the role of proof goal retention in the overall
optimization-based dependent typing architecture. The recurring theme is that
Pythonic source programs expose many plausible semantic candidates, but the
system must decide which of them should become trusted core obligations, which
should be surfaced as suggested annotations, which should remain as proof
goals, and which should be discarded as too unstable, too expensive, or too
poorly supported. This makes the architecture simultaneously a typing system,
a semantic extraction engine, a theory-integration framework, and an explicit
optimization procedure over candidate dependent information objects.

This supplementary note elaborates the role of proof goal retention in the overall
optimization-based dependent typing architecture. The recurring theme is that
Pythonic source programs expose many plausible semantic candidates, but the
system must decide which of them should become trusted core obligations, which
should be surfaced as suggested annotations, which should remain as proof
goals, and which should be discarded as too unstable, too expensive, or too
poorly supported. This makes the architecture simultaneously a typing system,
a semantic extraction engine, a theory-integration framework, and an explicit
optimization procedure over candidate dependent information objects.

This supplementary note elaborates the role of proof goal retention in the overall
optimization-based dependent typing architecture. The recurring theme is that
Pythonic source programs expose many plausible semantic candidates, but the
system must decide which of them should become trusted core obligations, which
should be surfaced as suggested annotations, which should remain as proof
goals, and which should be discarded as too unstable, too expensive, or too
poorly supported. This makes the architecture simultaneously a typing system,
a semantic extraction engine, a theory-integration framework, and an explicit
optimization procedure over candidate dependent information objects.

This supplementary note elaborates the role of proof goal retention in the overall
optimization-based dependent typing architecture. The recurring theme is that
Pythonic source programs expose many plausible semantic candidates, but the
system must decide which of them should become trusted core obligations, which
should be surfaced as suggested annotations, which should remain as proof
goals, and which should be discarded as too unstable, too expensive, or too
poorly supported. This makes the architecture simultaneously a typing system,
a semantic extraction engine, a theory-integration framework, and an explicit
optimization procedure over candidate dependent information objects.

\subsection{Supplement: annotation compression}
This supplementary note elaborates the role of annotation compression in the overall
optimization-based dependent typing architecture. The recurring theme is that
Pythonic source programs expose many plausible semantic candidates, but the
system must decide which of them should become trusted core obligations, which
should be surfaced as suggested annotations, which should remain as proof
goals, and which should be discarded as too unstable, too expensive, or too
poorly supported. This makes the architecture simultaneously a typing system,
a semantic extraction engine, a theory-integration framework, and an explicit
optimization procedure over candidate dependent information objects.

This supplementary note elaborates the role of annotation compression in the overall
optimization-based dependent typing architecture. The recurring theme is that
Pythonic source programs expose many plausible semantic candidates, but the
system must decide which of them should become trusted core obligations, which
should be surfaced as suggested annotations, which should remain as proof
goals, and which should be discarded as too unstable, too expensive, or too
poorly supported. This makes the architecture simultaneously a typing system,
a semantic extraction engine, a theory-integration framework, and an explicit
optimization procedure over candidate dependent information objects.

This supplementary note elaborates the role of annotation compression in the overall
optimization-based dependent typing architecture. The recurring theme is that
Pythonic source programs expose many plausible semantic candidates, but the
system must decide which of them should become trusted core obligations, which
should be surfaced as suggested annotations, which should remain as proof
goals, and which should be discarded as too unstable, too expensive, or too
poorly supported. This makes the architecture simultaneously a typing system,
a semantic extraction engine, a theory-integration framework, and an explicit
optimization procedure over candidate dependent information objects.

This supplementary note elaborates the role of annotation compression in the overall
optimization-based dependent typing architecture. The recurring theme is that
Pythonic source programs expose many plausible semantic candidates, but the
system must decide which of them should become trusted core obligations, which
should be surfaced as suggested annotations, which should remain as proof
goals, and which should be discarded as too unstable, too expensive, or too
poorly supported. This makes the architecture simultaneously a typing system,
a semantic extraction engine, a theory-integration framework, and an explicit
optimization procedure over candidate dependent information objects.

This supplementary note elaborates the role of annotation compression in the overall
optimization-based dependent typing architecture. The recurring theme is that
Pythonic source programs expose many plausible semantic candidates, but the
system must decide which of them should become trusted core obligations, which
should be surfaced as suggested annotations, which should remain as proof
goals, and which should be discarded as too unstable, too expensive, or too
poorly supported. This makes the architecture simultaneously a typing system,
a semantic extraction engine, a theory-integration framework, and an explicit
optimization procedure over candidate dependent information objects.

This supplementary note elaborates the role of annotation compression in the overall
optimization-based dependent typing architecture. The recurring theme is that
Pythonic source programs expose many plausible semantic candidates, but the
system must decide which of them should become trusted core obligations, which
should be surfaced as suggested annotations, which should remain as proof
goals, and which should be discarded as too unstable, too expensive, or too
poorly supported. This makes the architecture simultaneously a typing system,
a semantic extraction engine, a theory-integration framework, and an explicit
optimization procedure over candidate dependent information objects.

\subsection{Supplement: protocol entailment}
This supplementary note elaborates the role of protocol entailment in the overall
optimization-based dependent typing architecture. The recurring theme is that
Pythonic source programs expose many plausible semantic candidates, but the
system must decide which of them should become trusted core obligations, which
should be surfaced as suggested annotations, which should remain as proof
goals, and which should be discarded as too unstable, too expensive, or too
poorly supported. This makes the architecture simultaneously a typing system,
a semantic extraction engine, a theory-integration framework, and an explicit
optimization procedure over candidate dependent information objects.

This supplementary note elaborates the role of protocol entailment in the overall
optimization-based dependent typing architecture. The recurring theme is that
Pythonic source programs expose many plausible semantic candidates, but the
system must decide which of them should become trusted core obligations, which
should be surfaced as suggested annotations, which should remain as proof
goals, and which should be discarded as too unstable, too expensive, or too
poorly supported. This makes the architecture simultaneously a typing system,
a semantic extraction engine, a theory-integration framework, and an explicit
optimization procedure over candidate dependent information objects.

This supplementary note elaborates the role of protocol entailment in the overall
optimization-based dependent typing architecture. The recurring theme is that
Pythonic source programs expose many plausible semantic candidates, but the
system must decide which of them should become trusted core obligations, which
should be surfaced as suggested annotations, which should remain as proof
goals, and which should be discarded as too unstable, too expensive, or too
poorly supported. This makes the architecture simultaneously a typing system,
a semantic extraction engine, a theory-integration framework, and an explicit
optimization procedure over candidate dependent information objects.

This supplementary note elaborates the role of protocol entailment in the overall
optimization-based dependent typing architecture. The recurring theme is that
Pythonic source programs expose many plausible semantic candidates, but the
system must decide which of them should become trusted core obligations, which
should be surfaced as suggested annotations, which should remain as proof
goals, and which should be discarded as too unstable, too expensive, or too
poorly supported. This makes the architecture simultaneously a typing system,
a semantic extraction engine, a theory-integration framework, and an explicit
optimization procedure over candidate dependent information objects.

This supplementary note elaborates the role of protocol entailment in the overall
optimization-based dependent typing architecture. The recurring theme is that
Pythonic source programs expose many plausible semantic candidates, but the
system must decide which of them should become trusted core obligations, which
should be surfaced as suggested annotations, which should remain as proof
goals, and which should be discarded as too unstable, too expensive, or too
poorly supported. This makes the architecture simultaneously a typing system,
a semantic extraction engine, a theory-integration framework, and an explicit
optimization procedure over candidate dependent information objects.

This supplementary note elaborates the role of protocol entailment in the overall
optimization-based dependent typing architecture. The recurring theme is that
Pythonic source programs expose many plausible semantic candidates, but the
system must decide which of them should become trusted core obligations, which
should be surfaced as suggested annotations, which should remain as proof
goals, and which should be discarded as too unstable, too expensive, or too
poorly supported. This makes the architecture simultaneously a typing system,
a semantic extraction engine, a theory-integration framework, and an explicit
optimization procedure over candidate dependent information objects.

\subsection{Supplement: heap alias control}
This supplementary note elaborates the role of heap alias control in the overall
optimization-based dependent typing architecture. The recurring theme is that
Pythonic source programs expose many plausible semantic candidates, but the
system must decide which of them should become trusted core obligations, which
should be surfaced as suggested annotations, which should remain as proof
goals, and which should be discarded as too unstable, too expensive, or too
poorly supported. This makes the architecture simultaneously a typing system,
a semantic extraction engine, a theory-integration framework, and an explicit
optimization procedure over candidate dependent information objects.

This supplementary note elaborates the role of heap alias control in the overall
optimization-based dependent typing architecture. The recurring theme is that
Pythonic source programs expose many plausible semantic candidates, but the
system must decide which of them should become trusted core obligations, which
should be surfaced as suggested annotations, which should remain as proof
goals, and which should be discarded as too unstable, too expensive, or too
poorly supported. This makes the architecture simultaneously a typing system,
a semantic extraction engine, a theory-integration framework, and an explicit
optimization procedure over candidate dependent information objects.

This supplementary note elaborates the role of heap alias control in the overall
optimization-based dependent typing architecture. The recurring theme is that
Pythonic source programs expose many plausible semantic candidates, but the
system must decide which of them should become trusted core obligations, which
should be surfaced as suggested annotations, which should remain as proof
goals, and which should be discarded as too unstable, too expensive, or too
poorly supported. This makes the architecture simultaneously a typing system,
a semantic extraction engine, a theory-integration framework, and an explicit
optimization procedure over candidate dependent information objects.

This supplementary note elaborates the role of heap alias control in the overall
optimization-based dependent typing architecture. The recurring theme is that
Pythonic source programs expose many plausible semantic candidates, but the
system must decide which of them should become trusted core obligations, which
should be surfaced as suggested annotations, which should remain as proof
goals, and which should be discarded as too unstable, too expensive, or too
poorly supported. This makes the architecture simultaneously a typing system,
a semantic extraction engine, a theory-integration framework, and an explicit
optimization procedure over candidate dependent information objects.

This supplementary note elaborates the role of heap alias control in the overall
optimization-based dependent typing architecture. The recurring theme is that
Pythonic source programs expose many plausible semantic candidates, but the
system must decide which of them should become trusted core obligations, which
should be surfaced as suggested annotations, which should remain as proof
goals, and which should be discarded as too unstable, too expensive, or too
poorly supported. This makes the architecture simultaneously a typing system,
a semantic extraction engine, a theory-integration framework, and an explicit
optimization procedure over candidate dependent information objects.

This supplementary note elaborates the role of heap alias control in the overall
optimization-based dependent typing architecture. The recurring theme is that
Pythonic source programs expose many plausible semantic candidates, but the
system must decide which of them should become trusted core obligations, which
should be surfaced as suggested annotations, which should remain as proof
goals, and which should be discarded as too unstable, too expensive, or too
poorly supported. This makes the architecture simultaneously a typing system,
a semantic extraction engine, a theory-integration framework, and an explicit
optimization procedure over candidate dependent information objects.

\subsection{Supplement: PyTorch view semantics}
This supplementary note elaborates the role of PyTorch view semantics in the overall
optimization-based dependent typing architecture. The recurring theme is that
Pythonic source programs expose many plausible semantic candidates, but the
system must decide which of them should become trusted core obligations, which
should be surfaced as suggested annotations, which should remain as proof
goals, and which should be discarded as too unstable, too expensive, or too
poorly supported. This makes the architecture simultaneously a typing system,
a semantic extraction engine, a theory-integration framework, and an explicit
optimization procedure over candidate dependent information objects.

This supplementary note elaborates the role of PyTorch view semantics in the overall
optimization-based dependent typing architecture. The recurring theme is that
Pythonic source programs expose many plausible semantic candidates, but the
system must decide which of them should become trusted core obligations, which
should be surfaced as suggested annotations, which should remain as proof
goals, and which should be discarded as too unstable, too expensive, or too
poorly supported. This makes the architecture simultaneously a typing system,
a semantic extraction engine, a theory-integration framework, and an explicit
optimization procedure over candidate dependent information objects.

This supplementary note elaborates the role of PyTorch view semantics in the overall
optimization-based dependent typing architecture. The recurring theme is that
Pythonic source programs expose many plausible semantic candidates, but the
system must decide which of them should become trusted core obligations, which
should be surfaced as suggested annotations, which should remain as proof
goals, and which should be discarded as too unstable, too expensive, or too
poorly supported. This makes the architecture simultaneously a typing system,
a semantic extraction engine, a theory-integration framework, and an explicit
optimization procedure over candidate dependent information objects.

This supplementary note elaborates the role of PyTorch view semantics in the overall
optimization-based dependent typing architecture. The recurring theme is that
Pythonic source programs expose many plausible semantic candidates, but the
system must decide which of them should become trusted core obligations, which
should be surfaced as suggested annotations, which should remain as proof
goals, and which should be discarded as too unstable, too expensive, or too
poorly supported. This makes the architecture simultaneously a typing system,
a semantic extraction engine, a theory-integration framework, and an explicit
optimization procedure over candidate dependent information objects.

This supplementary note elaborates the role of PyTorch view semantics in the overall
optimization-based dependent typing architecture. The recurring theme is that
Pythonic source programs expose many plausible semantic candidates, but the
system must decide which of them should become trusted core obligations, which
should be surfaced as suggested annotations, which should remain as proof
goals, and which should be discarded as too unstable, too expensive, or too
poorly supported. This makes the architecture simultaneously a typing system,
a semantic extraction engine, a theory-integration framework, and an explicit
optimization procedure over candidate dependent information objects.

This supplementary note elaborates the role of PyTorch view semantics in the overall
optimization-based dependent typing architecture. The recurring theme is that
Pythonic source programs expose many plausible semantic candidates, but the
system must decide which of them should become trusted core obligations, which
should be surfaced as suggested annotations, which should remain as proof
goals, and which should be discarded as too unstable, too expensive, or too
poorly supported. This makes the architecture simultaneously a typing system,
a semantic extraction engine, a theory-integration framework, and an explicit
optimization procedure over candidate dependent information objects.

\subsection{Supplement: tensor sorting obligations}
This supplementary note elaborates the role of tensor sorting obligations in the overall
optimization-based dependent typing architecture. The recurring theme is that
Pythonic source programs expose many plausible semantic candidates, but the
system must decide which of them should become trusted core obligations, which
should be surfaced as suggested annotations, which should remain as proof
goals, and which should be discarded as too unstable, too expensive, or too
poorly supported. This makes the architecture simultaneously a typing system,
a semantic extraction engine, a theory-integration framework, and an explicit
optimization procedure over candidate dependent information objects.

This supplementary note elaborates the role of tensor sorting obligations in the overall
optimization-based dependent typing architecture. The recurring theme is that
Pythonic source programs expose many plausible semantic candidates, but the
system must decide which of them should become trusted core obligations, which
should be surfaced as suggested annotations, which should remain as proof
goals, and which should be discarded as too unstable, too expensive, or too
poorly supported. This makes the architecture simultaneously a typing system,
a semantic extraction engine, a theory-integration framework, and an explicit
optimization procedure over candidate dependent information objects.

This supplementary note elaborates the role of tensor sorting obligations in the overall
optimization-based dependent typing architecture. The recurring theme is that
Pythonic source programs expose many plausible semantic candidates, but the
system must decide which of them should become trusted core obligations, which
should be surfaced as suggested annotations, which should remain as proof
goals, and which should be discarded as too unstable, too expensive, or too
poorly supported. This makes the architecture simultaneously a typing system,
a semantic extraction engine, a theory-integration framework, and an explicit
optimization procedure over candidate dependent information objects.

This supplementary note elaborates the role of tensor sorting obligations in the overall
optimization-based dependent typing architecture. The recurring theme is that
Pythonic source programs expose many plausible semantic candidates, but the
system must decide which of them should become trusted core obligations, which
should be surfaced as suggested annotations, which should remain as proof
goals, and which should be discarded as too unstable, too expensive, or too
poorly supported. This makes the architecture simultaneously a typing system,
a semantic extraction engine, a theory-integration framework, and an explicit
optimization procedure over candidate dependent information objects.

This supplementary note elaborates the role of tensor sorting obligations in the overall
optimization-based dependent typing architecture. The recurring theme is that
Pythonic source programs expose many plausible semantic candidates, but the
system must decide which of them should become trusted core obligations, which
should be surfaced as suggested annotations, which should remain as proof
goals, and which should be discarded as too unstable, too expensive, or too
poorly supported. This makes the architecture simultaneously a typing system,
a semantic extraction engine, a theory-integration framework, and an explicit
optimization procedure over candidate dependent information objects.

This supplementary note elaborates the role of tensor sorting obligations in the overall
optimization-based dependent typing architecture. The recurring theme is that
Pythonic source programs expose many plausible semantic candidates, but the
system must decide which of them should become trusted core obligations, which
should be surfaced as suggested annotations, which should remain as proof
goals, and which should be discarded as too unstable, too expensive, or too
poorly supported. This makes the architecture simultaneously a typing system,
a semantic extraction engine, a theory-integration framework, and an explicit
optimization procedure over candidate dependent information objects.

\subsection{Supplement: qualifier bank management}
This supplementary note elaborates the role of qualifier bank management in the overall
optimization-based dependent typing architecture. The recurring theme is that
Pythonic source programs expose many plausible semantic candidates, but the
system must decide which of them should become trusted core obligations, which
should be surfaced as suggested annotations, which should remain as proof
goals, and which should be discarded as too unstable, too expensive, or too
poorly supported. This makes the architecture simultaneously a typing system,
a semantic extraction engine, a theory-integration framework, and an explicit
optimization procedure over candidate dependent information objects.

This supplementary note elaborates the role of qualifier bank management in the overall
optimization-based dependent typing architecture. The recurring theme is that
Pythonic source programs expose many plausible semantic candidates, but the
system must decide which of them should become trusted core obligations, which
should be surfaced as suggested annotations, which should remain as proof
goals, and which should be discarded as too unstable, too expensive, or too
poorly supported. This makes the architecture simultaneously a typing system,
a semantic extraction engine, a theory-integration framework, and an explicit
optimization procedure over candidate dependent information objects.

This supplementary note elaborates the role of qualifier bank management in the overall
optimization-based dependent typing architecture. The recurring theme is that
Pythonic source programs expose many plausible semantic candidates, but the
system must decide which of them should become trusted core obligations, which
should be surfaced as suggested annotations, which should remain as proof
goals, and which should be discarded as too unstable, too expensive, or too
poorly supported. This makes the architecture simultaneously a typing system,
a semantic extraction engine, a theory-integration framework, and an explicit
optimization procedure over candidate dependent information objects.

This supplementary note elaborates the role of qualifier bank management in the overall
optimization-based dependent typing architecture. The recurring theme is that
Pythonic source programs expose many plausible semantic candidates, but the
system must decide which of them should become trusted core obligations, which
should be surfaced as suggested annotations, which should remain as proof
goals, and which should be discarded as too unstable, too expensive, or too
poorly supported. This makes the architecture simultaneously a typing system,
a semantic extraction engine, a theory-integration framework, and an explicit
optimization procedure over candidate dependent information objects.

This supplementary note elaborates the role of qualifier bank management in the overall
optimization-based dependent typing architecture. The recurring theme is that
Pythonic source programs expose many plausible semantic candidates, but the
system must decide which of them should become trusted core obligations, which
should be surfaced as suggested annotations, which should remain as proof
goals, and which should be discarded as too unstable, too expensive, or too
poorly supported. This makes the architecture simultaneously a typing system,
a semantic extraction engine, a theory-integration framework, and an explicit
optimization procedure over candidate dependent information objects.

This supplementary note elaborates the role of qualifier bank management in the overall
optimization-based dependent typing architecture. The recurring theme is that
Pythonic source programs expose many plausible semantic candidates, but the
system must decide which of them should become trusted core obligations, which
should be surfaced as suggested annotations, which should remain as proof
goals, and which should be discarded as too unstable, too expensive, or too
poorly supported. This makes the architecture simultaneously a typing system,
a semantic extraction engine, a theory-integration framework, and an explicit
optimization procedure over candidate dependent information objects.

\subsection{Supplement: library plugin APIs}
This supplementary note elaborates the role of library plugin APIs in the overall
optimization-based dependent typing architecture. The recurring theme is that
Pythonic source programs expose many plausible semantic candidates, but the
system must decide which of them should become trusted core obligations, which
should be surfaced as suggested annotations, which should remain as proof
goals, and which should be discarded as too unstable, too expensive, or too
poorly supported. This makes the architecture simultaneously a typing system,
a semantic extraction engine, a theory-integration framework, and an explicit
optimization procedure over candidate dependent information objects.

This supplementary note elaborates the role of library plugin APIs in the overall
optimization-based dependent typing architecture. The recurring theme is that
Pythonic source programs expose many plausible semantic candidates, but the
system must decide which of them should become trusted core obligations, which
should be surfaced as suggested annotations, which should remain as proof
goals, and which should be discarded as too unstable, too expensive, or too
poorly supported. This makes the architecture simultaneously a typing system,
a semantic extraction engine, a theory-integration framework, and an explicit
optimization procedure over candidate dependent information objects.

This supplementary note elaborates the role of library plugin APIs in the overall
optimization-based dependent typing architecture. The recurring theme is that
Pythonic source programs expose many plausible semantic candidates, but the
system must decide which of them should become trusted core obligations, which
should be surfaced as suggested annotations, which should remain as proof
goals, and which should be discarded as too unstable, too expensive, or too
poorly supported. This makes the architecture simultaneously a typing system,
a semantic extraction engine, a theory-integration framework, and an explicit
optimization procedure over candidate dependent information objects.

This supplementary note elaborates the role of library plugin APIs in the overall
optimization-based dependent typing architecture. The recurring theme is that
Pythonic source programs expose many plausible semantic candidates, but the
system must decide which of them should become trusted core obligations, which
should be surfaced as suggested annotations, which should remain as proof
goals, and which should be discarded as too unstable, too expensive, or too
poorly supported. This makes the architecture simultaneously a typing system,
a semantic extraction engine, a theory-integration framework, and an explicit
optimization procedure over candidate dependent information objects.

This supplementary note elaborates the role of library plugin APIs in the overall
optimization-based dependent typing architecture. The recurring theme is that
Pythonic source programs expose many plausible semantic candidates, but the
system must decide which of them should become trusted core obligations, which
should be surfaced as suggested annotations, which should remain as proof
goals, and which should be discarded as too unstable, too expensive, or too
poorly supported. This makes the architecture simultaneously a typing system,
a semantic extraction engine, a theory-integration framework, and an explicit
optimization procedure over candidate dependent information objects.

This supplementary note elaborates the role of library plugin APIs in the overall
optimization-based dependent typing architecture. The recurring theme is that
Pythonic source programs expose many plausible semantic candidates, but the
system must decide which of them should become trusted core obligations, which
should be surfaced as suggested annotations, which should remain as proof
goals, and which should be discarded as too unstable, too expensive, or too
poorly supported. This makes the architecture simultaneously a typing system,
a semantic extraction engine, a theory-integration framework, and an explicit
optimization procedure over candidate dependent information objects.

\subsection{Supplement: LLM-authored guarantee ingestion}
This supplementary note elaborates the role of LLM-authored guarantee ingestion in the overall
optimization-based dependent typing architecture. The recurring theme is that
Pythonic source programs expose many plausible semantic candidates, but the
system must decide which of them should become trusted core obligations, which
should be surfaced as suggested annotations, which should remain as proof
goals, and which should be discarded as too unstable, too expensive, or too
poorly supported. This makes the architecture simultaneously a typing system,
a semantic extraction engine, a theory-integration framework, and an explicit
optimization procedure over candidate dependent information objects.

This supplementary note elaborates the role of LLM-authored guarantee ingestion in the overall
optimization-based dependent typing architecture. The recurring theme is that
Pythonic source programs expose many plausible semantic candidates, but the
system must decide which of them should become trusted core obligations, which
should be surfaced as suggested annotations, which should remain as proof
goals, and which should be discarded as too unstable, too expensive, or too
poorly supported. This makes the architecture simultaneously a typing system,
a semantic extraction engine, a theory-integration framework, and an explicit
optimization procedure over candidate dependent information objects.

This supplementary note elaborates the role of LLM-authored guarantee ingestion in the overall
optimization-based dependent typing architecture. The recurring theme is that
Pythonic source programs expose many plausible semantic candidates, but the
system must decide which of them should become trusted core obligations, which
should be surfaced as suggested annotations, which should remain as proof
goals, and which should be discarded as too unstable, too expensive, or too
poorly supported. This makes the architecture simultaneously a typing system,
a semantic extraction engine, a theory-integration framework, and an explicit
optimization procedure over candidate dependent information objects.

This supplementary note elaborates the role of LLM-authored guarantee ingestion in the overall
optimization-based dependent typing architecture. The recurring theme is that
Pythonic source programs expose many plausible semantic candidates, but the
system must decide which of them should become trusted core obligations, which
should be surfaced as suggested annotations, which should remain as proof
goals, and which should be discarded as too unstable, too expensive, or too
poorly supported. This makes the architecture simultaneously a typing system,
a semantic extraction engine, a theory-integration framework, and an explicit
optimization procedure over candidate dependent information objects.

This supplementary note elaborates the role of LLM-authored guarantee ingestion in the overall
optimization-based dependent typing architecture. The recurring theme is that
Pythonic source programs expose many plausible semantic candidates, but the
system must decide which of them should become trusted core obligations, which
should be surfaced as suggested annotations, which should remain as proof
goals, and which should be discarded as too unstable, too expensive, or too
poorly supported. This makes the architecture simultaneously a typing system,
a semantic extraction engine, a theory-integration framework, and an explicit
optimization procedure over candidate dependent information objects.

This supplementary note elaborates the role of LLM-authored guarantee ingestion in the overall
optimization-based dependent typing architecture. The recurring theme is that
Pythonic source programs expose many plausible semantic candidates, but the
system must decide which of them should become trusted core obligations, which
should be surfaced as suggested annotations, which should remain as proof
goals, and which should be discarded as too unstable, too expensive, or too
poorly supported. This makes the architecture simultaneously a typing system,
a semantic extraction engine, a theory-integration framework, and an explicit
optimization procedure over candidate dependent information objects.

\subsection{Supplement: counterexample-guided qualifier repair}
This supplementary note elaborates the role of counterexample-guided qualifier repair in the overall
optimization-based dependent typing architecture. The recurring theme is that
Pythonic source programs expose many plausible semantic candidates, but the
system must decide which of them should become trusted core obligations, which
should be surfaced as suggested annotations, which should remain as proof
goals, and which should be discarded as too unstable, too expensive, or too
poorly supported. This makes the architecture simultaneously a typing system,
a semantic extraction engine, a theory-integration framework, and an explicit
optimization procedure over candidate dependent information objects.

This supplementary note elaborates the role of counterexample-guided qualifier repair in the overall
optimization-based dependent typing architecture. The recurring theme is that
Pythonic source programs expose many plausible semantic candidates, but the
system must decide which of them should become trusted core obligations, which
should be surfaced as suggested annotations, which should remain as proof
goals, and which should be discarded as too unstable, too expensive, or too
poorly supported. This makes the architecture simultaneously a typing system,
a semantic extraction engine, a theory-integration framework, and an explicit
optimization procedure over candidate dependent information objects.

This supplementary note elaborates the role of counterexample-guided qualifier repair in the overall
optimization-based dependent typing architecture. The recurring theme is that
Pythonic source programs expose many plausible semantic candidates, but the
system must decide which of them should become trusted core obligations, which
should be surfaced as suggested annotations, which should remain as proof
goals, and which should be discarded as too unstable, too expensive, or too
poorly supported. This makes the architecture simultaneously a typing system,
a semantic extraction engine, a theory-integration framework, and an explicit
optimization procedure over candidate dependent information objects.

This supplementary note elaborates the role of counterexample-guided qualifier repair in the overall
optimization-based dependent typing architecture. The recurring theme is that
Pythonic source programs expose many plausible semantic candidates, but the
system must decide which of them should become trusted core obligations, which
should be surfaced as suggested annotations, which should remain as proof
goals, and which should be discarded as too unstable, too expensive, or too
poorly supported. This makes the architecture simultaneously a typing system,
a semantic extraction engine, a theory-integration framework, and an explicit
optimization procedure over candidate dependent information objects.

This supplementary note elaborates the role of counterexample-guided qualifier repair in the overall
optimization-based dependent typing architecture. The recurring theme is that
Pythonic source programs expose many plausible semantic candidates, but the
system must decide which of them should become trusted core obligations, which
should be surfaced as suggested annotations, which should remain as proof
goals, and which should be discarded as too unstable, too expensive, or too
poorly supported. This makes the architecture simultaneously a typing system,
a semantic extraction engine, a theory-integration framework, and an explicit
optimization procedure over candidate dependent information objects.

This supplementary note elaborates the role of counterexample-guided qualifier repair in the overall
optimization-based dependent typing architecture. The recurring theme is that
Pythonic source programs expose many plausible semantic candidates, but the
system must decide which of them should become trusted core obligations, which
should be surfaced as suggested annotations, which should remain as proof
goals, and which should be discarded as too unstable, too expensive, or too
poorly supported. This makes the architecture simultaneously a typing system,
a semantic extraction engine, a theory-integration framework, and an explicit
optimization procedure over candidate dependent information objects.

\subsection{Supplement: subtyping under theory refinement}
This supplementary note elaborates the role of subtyping under theory refinement in the overall
optimization-based dependent typing architecture. The recurring theme is that
Pythonic source programs expose many plausible semantic candidates, but the
system must decide which of them should become trusted core obligations, which
should be surfaced as suggested annotations, which should remain as proof
goals, and which should be discarded as too unstable, too expensive, or too
poorly supported. This makes the architecture simultaneously a typing system,
a semantic extraction engine, a theory-integration framework, and an explicit
optimization procedure over candidate dependent information objects.

This supplementary note elaborates the role of subtyping under theory refinement in the overall
optimization-based dependent typing architecture. The recurring theme is that
Pythonic source programs expose many plausible semantic candidates, but the
system must decide which of them should become trusted core obligations, which
should be surfaced as suggested annotations, which should remain as proof
goals, and which should be discarded as too unstable, too expensive, or too
poorly supported. This makes the architecture simultaneously a typing system,
a semantic extraction engine, a theory-integration framework, and an explicit
optimization procedure over candidate dependent information objects.

This supplementary note elaborates the role of subtyping under theory refinement in the overall
optimization-based dependent typing architecture. The recurring theme is that
Pythonic source programs expose many plausible semantic candidates, but the
system must decide which of them should become trusted core obligations, which
should be surfaced as suggested annotations, which should remain as proof
goals, and which should be discarded as too unstable, too expensive, or too
poorly supported. This makes the architecture simultaneously a typing system,
a semantic extraction engine, a theory-integration framework, and an explicit
optimization procedure over candidate dependent information objects.

This supplementary note elaborates the role of subtyping under theory refinement in the overall
optimization-based dependent typing architecture. The recurring theme is that
Pythonic source programs expose many plausible semantic candidates, but the
system must decide which of them should become trusted core obligations, which
should be surfaced as suggested annotations, which should remain as proof
goals, and which should be discarded as too unstable, too expensive, or too
poorly supported. This makes the architecture simultaneously a typing system,
a semantic extraction engine, a theory-integration framework, and an explicit
optimization procedure over candidate dependent information objects.

This supplementary note elaborates the role of subtyping under theory refinement in the overall
optimization-based dependent typing architecture. The recurring theme is that
Pythonic source programs expose many plausible semantic candidates, but the
system must decide which of them should become trusted core obligations, which
should be surfaced as suggested annotations, which should remain as proof
goals, and which should be discarded as too unstable, too expensive, or too
poorly supported. This makes the architecture simultaneously a typing system,
a semantic extraction engine, a theory-integration framework, and an explicit
optimization procedure over candidate dependent information objects.

This supplementary note elaborates the role of subtyping under theory refinement in the overall
optimization-based dependent typing architecture. The recurring theme is that
Pythonic source programs expose many plausible semantic candidates, but the
system must decide which of them should become trusted core obligations, which
should be surfaced as suggested annotations, which should remain as proof
goals, and which should be discarded as too unstable, too expensive, or too
poorly supported. This makes the architecture simultaneously a typing system,
a semantic extraction engine, a theory-integration framework, and an explicit
optimization procedure over candidate dependent information objects.

\subsection{Supplement: constructor invariants}
This supplementary note elaborates the role of constructor invariants in the overall
optimization-based dependent typing architecture. The recurring theme is that
Pythonic source programs expose many plausible semantic candidates, but the
system must decide which of them should become trusted core obligations, which
should be surfaced as suggested annotations, which should remain as proof
goals, and which should be discarded as too unstable, too expensive, or too
poorly supported. This makes the architecture simultaneously a typing system,
a semantic extraction engine, a theory-integration framework, and an explicit
optimization procedure over candidate dependent information objects.

This supplementary note elaborates the role of constructor invariants in the overall
optimization-based dependent typing architecture. The recurring theme is that
Pythonic source programs expose many plausible semantic candidates, but the
system must decide which of them should become trusted core obligations, which
should be surfaced as suggested annotations, which should remain as proof
goals, and which should be discarded as too unstable, too expensive, or too
poorly supported. This makes the architecture simultaneously a typing system,
a semantic extraction engine, a theory-integration framework, and an explicit
optimization procedure over candidate dependent information objects.

This supplementary note elaborates the role of constructor invariants in the overall
optimization-based dependent typing architecture. The recurring theme is that
Pythonic source programs expose many plausible semantic candidates, but the
system must decide which of them should become trusted core obligations, which
should be surfaced as suggested annotations, which should remain as proof
goals, and which should be discarded as too unstable, too expensive, or too
poorly supported. This makes the architecture simultaneously a typing system,
a semantic extraction engine, a theory-integration framework, and an explicit
optimization procedure over candidate dependent information objects.

This supplementary note elaborates the role of constructor invariants in the overall
optimization-based dependent typing architecture. The recurring theme is that
Pythonic source programs expose many plausible semantic candidates, but the
system must decide which of them should become trusted core obligations, which
should be surfaced as suggested annotations, which should remain as proof
goals, and which should be discarded as too unstable, too expensive, or too
poorly supported. This makes the architecture simultaneously a typing system,
a semantic extraction engine, a theory-integration framework, and an explicit
optimization procedure over candidate dependent information objects.

This supplementary note elaborates the role of constructor invariants in the overall
optimization-based dependent typing architecture. The recurring theme is that
Pythonic source programs expose many plausible semantic candidates, but the
system must decide which of them should become trusted core obligations, which
should be surfaced as suggested annotations, which should remain as proof
goals, and which should be discarded as too unstable, too expensive, or too
poorly supported. This makes the architecture simultaneously a typing system,
a semantic extraction engine, a theory-integration framework, and an explicit
optimization procedure over candidate dependent information objects.

This supplementary note elaborates the role of constructor invariants in the overall
optimization-based dependent typing architecture. The recurring theme is that
Pythonic source programs expose many plausible semantic candidates, but the
system must decide which of them should become trusted core obligations, which
should be surfaced as suggested annotations, which should remain as proof
goals, and which should be discarded as too unstable, too expensive, or too
poorly supported. This makes the architecture simultaneously a typing system,
a semantic extraction engine, a theory-integration framework, and an explicit
optimization procedure over candidate dependent information objects.

\subsection{Supplement: mutation-stability filtering}
This supplementary note elaborates the role of mutation-stability filtering in the overall
optimization-based dependent typing architecture. The recurring theme is that
Pythonic source programs expose many plausible semantic candidates, but the
system must decide which of them should become trusted core obligations, which
should be surfaced as suggested annotations, which should remain as proof
goals, and which should be discarded as too unstable, too expensive, or too
poorly supported. This makes the architecture simultaneously a typing system,
a semantic extraction engine, a theory-integration framework, and an explicit
optimization procedure over candidate dependent information objects.

This supplementary note elaborates the role of mutation-stability filtering in the overall
optimization-based dependent typing architecture. The recurring theme is that
Pythonic source programs expose many plausible semantic candidates, but the
system must decide which of them should become trusted core obligations, which
should be surfaced as suggested annotations, which should remain as proof
goals, and which should be discarded as too unstable, too expensive, or too
poorly supported. This makes the architecture simultaneously a typing system,
a semantic extraction engine, a theory-integration framework, and an explicit
optimization procedure over candidate dependent information objects.

This supplementary note elaborates the role of mutation-stability filtering in the overall
optimization-based dependent typing architecture. The recurring theme is that
Pythonic source programs expose many plausible semantic candidates, but the
system must decide which of them should become trusted core obligations, which
should be surfaced as suggested annotations, which should remain as proof
goals, and which should be discarded as too unstable, too expensive, or too
poorly supported. This makes the architecture simultaneously a typing system,
a semantic extraction engine, a theory-integration framework, and an explicit
optimization procedure over candidate dependent information objects.

This supplementary note elaborates the role of mutation-stability filtering in the overall
optimization-based dependent typing architecture. The recurring theme is that
Pythonic source programs expose many plausible semantic candidates, but the
system must decide which of them should become trusted core obligations, which
should be surfaced as suggested annotations, which should remain as proof
goals, and which should be discarded as too unstable, too expensive, or too
poorly supported. This makes the architecture simultaneously a typing system,
a semantic extraction engine, a theory-integration framework, and an explicit
optimization procedure over candidate dependent information objects.

This supplementary note elaborates the role of mutation-stability filtering in the overall
optimization-based dependent typing architecture. The recurring theme is that
Pythonic source programs expose many plausible semantic candidates, but the
system must decide which of them should become trusted core obligations, which
should be surfaced as suggested annotations, which should remain as proof
goals, and which should be discarded as too unstable, too expensive, or too
poorly supported. This makes the architecture simultaneously a typing system,
a semantic extraction engine, a theory-integration framework, and an explicit
optimization procedure over candidate dependent information objects.

This supplementary note elaborates the role of mutation-stability filtering in the overall
optimization-based dependent typing architecture. The recurring theme is that
Pythonic source programs expose many plausible semantic candidates, but the
system must decide which of them should become trusted core obligations, which
should be surfaced as suggested annotations, which should remain as proof
goals, and which should be discarded as too unstable, too expensive, or too
poorly supported. This makes the architecture simultaneously a typing system,
a semantic extraction engine, a theory-integration framework, and an explicit
optimization procedure over candidate dependent information objects.

\subsection{Supplement: batching and sequence models}
This supplementary note elaborates the role of batching and sequence models in the overall
optimization-based dependent typing architecture. The recurring theme is that
Pythonic source programs expose many plausible semantic candidates, but the
system must decide which of them should become trusted core obligations, which
should be surfaced as suggested annotations, which should remain as proof
goals, and which should be discarded as too unstable, too expensive, or too
poorly supported. This makes the architecture simultaneously a typing system,
a semantic extraction engine, a theory-integration framework, and an explicit
optimization procedure over candidate dependent information objects.

This supplementary note elaborates the role of batching and sequence models in the overall
optimization-based dependent typing architecture. The recurring theme is that
Pythonic source programs expose many plausible semantic candidates, but the
system must decide which of them should become trusted core obligations, which
should be surfaced as suggested annotations, which should remain as proof
goals, and which should be discarded as too unstable, too expensive, or too
poorly supported. This makes the architecture simultaneously a typing system,
a semantic extraction engine, a theory-integration framework, and an explicit
optimization procedure over candidate dependent information objects.

This supplementary note elaborates the role of batching and sequence models in the overall
optimization-based dependent typing architecture. The recurring theme is that
Pythonic source programs expose many plausible semantic candidates, but the
system must decide which of them should become trusted core obligations, which
should be surfaced as suggested annotations, which should remain as proof
goals, and which should be discarded as too unstable, too expensive, or too
poorly supported. This makes the architecture simultaneously a typing system,
a semantic extraction engine, a theory-integration framework, and an explicit
optimization procedure over candidate dependent information objects.

This supplementary note elaborates the role of batching and sequence models in the overall
optimization-based dependent typing architecture. The recurring theme is that
Pythonic source programs expose many plausible semantic candidates, but the
system must decide which of them should become trusted core obligations, which
should be surfaced as suggested annotations, which should remain as proof
goals, and which should be discarded as too unstable, too expensive, or too
poorly supported. This makes the architecture simultaneously a typing system,
a semantic extraction engine, a theory-integration framework, and an explicit
optimization procedure over candidate dependent information objects.

This supplementary note elaborates the role of batching and sequence models in the overall
optimization-based dependent typing architecture. The recurring theme is that
Pythonic source programs expose many plausible semantic candidates, but the
system must decide which of them should become trusted core obligations, which
should be surfaced as suggested annotations, which should remain as proof
goals, and which should be discarded as too unstable, too expensive, or too
poorly supported. This makes the architecture simultaneously a typing system,
a semantic extraction engine, a theory-integration framework, and an explicit
optimization procedure over candidate dependent information objects.

This supplementary note elaborates the role of batching and sequence models in the overall
optimization-based dependent typing architecture. The recurring theme is that
Pythonic source programs expose many plausible semantic candidates, but the
system must decide which of them should become trusted core obligations, which
should be surfaced as suggested annotations, which should remain as proof
goals, and which should be discarded as too unstable, too expensive, or too
poorly supported. This makes the architecture simultaneously a typing system,
a semantic extraction engine, a theory-integration framework, and an explicit
optimization procedure over candidate dependent information objects.

\subsection{Supplement: optimizer telemetry}
This supplementary note elaborates the role of optimizer telemetry in the overall
optimization-based dependent typing architecture. The recurring theme is that
Pythonic source programs expose many plausible semantic candidates, but the
system must decide which of them should become trusted core obligations, which
should be surfaced as suggested annotations, which should remain as proof
goals, and which should be discarded as too unstable, too expensive, or too
poorly supported. This makes the architecture simultaneously a typing system,
a semantic extraction engine, a theory-integration framework, and an explicit
optimization procedure over candidate dependent information objects.

This supplementary note elaborates the role of optimizer telemetry in the overall
optimization-based dependent typing architecture. The recurring theme is that
Pythonic source programs expose many plausible semantic candidates, but the
system must decide which of them should become trusted core obligations, which
should be surfaced as suggested annotations, which should remain as proof
goals, and which should be discarded as too unstable, too expensive, or too
poorly supported. This makes the architecture simultaneously a typing system,
a semantic extraction engine, a theory-integration framework, and an explicit
optimization procedure over candidate dependent information objects.

This supplementary note elaborates the role of optimizer telemetry in the overall
optimization-based dependent typing architecture. The recurring theme is that
Pythonic source programs expose many plausible semantic candidates, but the
system must decide which of them should become trusted core obligations, which
should be surfaced as suggested annotations, which should remain as proof
goals, and which should be discarded as too unstable, too expensive, or too
poorly supported. This makes the architecture simultaneously a typing system,
a semantic extraction engine, a theory-integration framework, and an explicit
optimization procedure over candidate dependent information objects.

This supplementary note elaborates the role of optimizer telemetry in the overall
optimization-based dependent typing architecture. The recurring theme is that
Pythonic source programs expose many plausible semantic candidates, but the
system must decide which of them should become trusted core obligations, which
should be surfaced as suggested annotations, which should remain as proof
goals, and which should be discarded as too unstable, too expensive, or too
poorly supported. This makes the architecture simultaneously a typing system,
a semantic extraction engine, a theory-integration framework, and an explicit
optimization procedure over candidate dependent information objects.

This supplementary note elaborates the role of optimizer telemetry in the overall
optimization-based dependent typing architecture. The recurring theme is that
Pythonic source programs expose many plausible semantic candidates, but the
system must decide which of them should become trusted core obligations, which
should be surfaced as suggested annotations, which should remain as proof
goals, and which should be discarded as too unstable, too expensive, or too
poorly supported. This makes the architecture simultaneously a typing system,
a semantic extraction engine, a theory-integration framework, and an explicit
optimization procedure over candidate dependent information objects.

This supplementary note elaborates the role of optimizer telemetry in the overall
optimization-based dependent typing architecture. The recurring theme is that
Pythonic source programs expose many plausible semantic candidates, but the
system must decide which of them should become trusted core obligations, which
should be surfaced as suggested annotations, which should remain as proof
goals, and which should be discarded as too unstable, too expensive, or too
poorly supported. This makes the architecture simultaneously a typing system,
a semantic extraction engine, a theory-integration framework, and an explicit
optimization procedure over candidate dependent information objects.

\subsection{Supplement: proof witness caching}
This supplementary note elaborates the role of proof witness caching in the overall
optimization-based dependent typing architecture. The recurring theme is that
Pythonic source programs expose many plausible semantic candidates, but the
system must decide which of them should become trusted core obligations, which
should be surfaced as suggested annotations, which should remain as proof
goals, and which should be discarded as too unstable, too expensive, or too
poorly supported. This makes the architecture simultaneously a typing system,
a semantic extraction engine, a theory-integration framework, and an explicit
optimization procedure over candidate dependent information objects.

This supplementary note elaborates the role of proof witness caching in the overall
optimization-based dependent typing architecture. The recurring theme is that
Pythonic source programs expose many plausible semantic candidates, but the
system must decide which of them should become trusted core obligations, which
should be surfaced as suggested annotations, which should remain as proof
goals, and which should be discarded as too unstable, too expensive, or too
poorly supported. This makes the architecture simultaneously a typing system,
a semantic extraction engine, a theory-integration framework, and an explicit
optimization procedure over candidate dependent information objects.

This supplementary note elaborates the role of proof witness caching in the overall
optimization-based dependent typing architecture. The recurring theme is that
Pythonic source programs expose many plausible semantic candidates, but the
system must decide which of them should become trusted core obligations, which
should be surfaced as suggested annotations, which should remain as proof
goals, and which should be discarded as too unstable, too expensive, or too
poorly supported. This makes the architecture simultaneously a typing system,
a semantic extraction engine, a theory-integration framework, and an explicit
optimization procedure over candidate dependent information objects.

This supplementary note elaborates the role of proof witness caching in the overall
optimization-based dependent typing architecture. The recurring theme is that
Pythonic source programs expose many plausible semantic candidates, but the
system must decide which of them should become trusted core obligations, which
should be surfaced as suggested annotations, which should remain as proof
goals, and which should be discarded as too unstable, too expensive, or too
poorly supported. This makes the architecture simultaneously a typing system,
a semantic extraction engine, a theory-integration framework, and an explicit
optimization procedure over candidate dependent information objects.

This supplementary note elaborates the role of proof witness caching in the overall
optimization-based dependent typing architecture. The recurring theme is that
Pythonic source programs expose many plausible semantic candidates, but the
system must decide which of them should become trusted core obligations, which
should be surfaced as suggested annotations, which should remain as proof
goals, and which should be discarded as too unstable, too expensive, or too
poorly supported. This makes the architecture simultaneously a typing system,
a semantic extraction engine, a theory-integration framework, and an explicit
optimization procedure over candidate dependent information objects.

This supplementary note elaborates the role of proof witness caching in the overall
optimization-based dependent typing architecture. The recurring theme is that
Pythonic source programs expose many plausible semantic candidates, but the
system must decide which of them should become trusted core obligations, which
should be surfaced as suggested annotations, which should remain as proof
goals, and which should be discarded as too unstable, too expensive, or too
poorly supported. This makes the architecture simultaneously a typing system,
a semantic extraction engine, a theory-integration framework, and an explicit
optimization procedure over candidate dependent information objects.

\subsection{Supplement: native-Python exception edges}
This supplementary note elaborates the role of native-Python exception edges in the overall
optimization-based dependent typing architecture. The recurring theme is that
Pythonic source programs expose many plausible semantic candidates, but the
system must decide which of them should become trusted core obligations, which
should be surfaced as suggested annotations, which should remain as proof
goals, and which should be discarded as too unstable, too expensive, or too
poorly supported. This makes the architecture simultaneously a typing system,
a semantic extraction engine, a theory-integration framework, and an explicit
optimization procedure over candidate dependent information objects.

This supplementary note elaborates the role of native-Python exception edges in the overall
optimization-based dependent typing architecture. The recurring theme is that
Pythonic source programs expose many plausible semantic candidates, but the
system must decide which of them should become trusted core obligations, which
should be surfaced as suggested annotations, which should remain as proof
goals, and which should be discarded as too unstable, too expensive, or too
poorly supported. This makes the architecture simultaneously a typing system,
a semantic extraction engine, a theory-integration framework, and an explicit
optimization procedure over candidate dependent information objects.

This supplementary note elaborates the role of native-Python exception edges in the overall
optimization-based dependent typing architecture. The recurring theme is that
Pythonic source programs expose many plausible semantic candidates, but the
system must decide which of them should become trusted core obligations, which
should be surfaced as suggested annotations, which should remain as proof
goals, and which should be discarded as too unstable, too expensive, or too
poorly supported. This makes the architecture simultaneously a typing system,
a semantic extraction engine, a theory-integration framework, and an explicit
optimization procedure over candidate dependent information objects.

This supplementary note elaborates the role of native-Python exception edges in the overall
optimization-based dependent typing architecture. The recurring theme is that
Pythonic source programs expose many plausible semantic candidates, but the
system must decide which of them should become trusted core obligations, which
should be surfaced as suggested annotations, which should remain as proof
goals, and which should be discarded as too unstable, too expensive, or too
poorly supported. This makes the architecture simultaneously a typing system,
a semantic extraction engine, a theory-integration framework, and an explicit
optimization procedure over candidate dependent information objects.

This supplementary note elaborates the role of native-Python exception edges in the overall
optimization-based dependent typing architecture. The recurring theme is that
Pythonic source programs expose many plausible semantic candidates, but the
system must decide which of them should become trusted core obligations, which
should be surfaced as suggested annotations, which should remain as proof
goals, and which should be discarded as too unstable, too expensive, or too
poorly supported. This makes the architecture simultaneously a typing system,
a semantic extraction engine, a theory-integration framework, and an explicit
optimization procedure over candidate dependent information objects.

This supplementary note elaborates the role of native-Python exception edges in the overall
optimization-based dependent typing architecture. The recurring theme is that
Pythonic source programs expose many plausible semantic candidates, but the
system must decide which of them should become trusted core obligations, which
should be surfaced as suggested annotations, which should remain as proof
goals, and which should be discarded as too unstable, too expensive, or too
poorly supported. This makes the architecture simultaneously a typing system,
a semantic extraction engine, a theory-integration framework, and an explicit
optimization procedure over candidate dependent information objects.

\subsection{Supplement: metaprogramming fallbacks}
This supplementary note elaborates the role of metaprogramming fallbacks in the overall
optimization-based dependent typing architecture. The recurring theme is that
Pythonic source programs expose many plausible semantic candidates, but the
system must decide which of them should become trusted core obligations, which
should be surfaced as suggested annotations, which should remain as proof
goals, and which should be discarded as too unstable, too expensive, or too
poorly supported. This makes the architecture simultaneously a typing system,
a semantic extraction engine, a theory-integration framework, and an explicit
optimization procedure over candidate dependent information objects.

This supplementary note elaborates the role of metaprogramming fallbacks in the overall
optimization-based dependent typing architecture. The recurring theme is that
Pythonic source programs expose many plausible semantic candidates, but the
system must decide which of them should become trusted core obligations, which
should be surfaced as suggested annotations, which should remain as proof
goals, and which should be discarded as too unstable, too expensive, or too
poorly supported. This makes the architecture simultaneously a typing system,
a semantic extraction engine, a theory-integration framework, and an explicit
optimization procedure over candidate dependent information objects.

This supplementary note elaborates the role of metaprogramming fallbacks in the overall
optimization-based dependent typing architecture. The recurring theme is that
Pythonic source programs expose many plausible semantic candidates, but the
system must decide which of them should become trusted core obligations, which
should be surfaced as suggested annotations, which should remain as proof
goals, and which should be discarded as too unstable, too expensive, or too
poorly supported. This makes the architecture simultaneously a typing system,
a semantic extraction engine, a theory-integration framework, and an explicit
optimization procedure over candidate dependent information objects.

This supplementary note elaborates the role of metaprogramming fallbacks in the overall
optimization-based dependent typing architecture. The recurring theme is that
Pythonic source programs expose many plausible semantic candidates, but the
system must decide which of them should become trusted core obligations, which
should be surfaced as suggested annotations, which should remain as proof
goals, and which should be discarded as too unstable, too expensive, or too
poorly supported. This makes the architecture simultaneously a typing system,
a semantic extraction engine, a theory-integration framework, and an explicit
optimization procedure over candidate dependent information objects.

This supplementary note elaborates the role of metaprogramming fallbacks in the overall
optimization-based dependent typing architecture. The recurring theme is that
Pythonic source programs expose many plausible semantic candidates, but the
system must decide which of them should become trusted core obligations, which
should be surfaced as suggested annotations, which should remain as proof
goals, and which should be discarded as too unstable, too expensive, or too
poorly supported. This makes the architecture simultaneously a typing system,
a semantic extraction engine, a theory-integration framework, and an explicit
optimization procedure over candidate dependent information objects.

This supplementary note elaborates the role of metaprogramming fallbacks in the overall
optimization-based dependent typing architecture. The recurring theme is that
Pythonic source programs expose many plausible semantic candidates, but the
system must decide which of them should become trusted core obligations, which
should be surfaced as suggested annotations, which should remain as proof
goals, and which should be discarded as too unstable, too expensive, or too
poorly supported. This makes the architecture simultaneously a typing system,
a semantic extraction engine, a theory-integration framework, and an explicit
optimization procedure over candidate dependent information objects.

\subsection{Supplement: trusted-kernel auditing}
This supplementary note elaborates the role of trusted-kernel auditing in the overall
optimization-based dependent typing architecture. The recurring theme is that
Pythonic source programs expose many plausible semantic candidates, but the
system must decide which of them should become trusted core obligations, which
should be surfaced as suggested annotations, which should remain as proof
goals, and which should be discarded as too unstable, too expensive, or too
poorly supported. This makes the architecture simultaneously a typing system,
a semantic extraction engine, a theory-integration framework, and an explicit
optimization procedure over candidate dependent information objects.

This supplementary note elaborates the role of trusted-kernel auditing in the overall
optimization-based dependent typing architecture. The recurring theme is that
Pythonic source programs expose many plausible semantic candidates, but the
system must decide which of them should become trusted core obligations, which
should be surfaced as suggested annotations, which should remain as proof
goals, and which should be discarded as too unstable, too expensive, or too
poorly supported. This makes the architecture simultaneously a typing system,
a semantic extraction engine, a theory-integration framework, and an explicit
optimization procedure over candidate dependent information objects.

This supplementary note elaborates the role of trusted-kernel auditing in the overall
optimization-based dependent typing architecture. The recurring theme is that
Pythonic source programs expose many plausible semantic candidates, but the
system must decide which of them should become trusted core obligations, which
should be surfaced as suggested annotations, which should remain as proof
goals, and which should be discarded as too unstable, too expensive, or too
poorly supported. This makes the architecture simultaneously a typing system,
a semantic extraction engine, a theory-integration framework, and an explicit
optimization procedure over candidate dependent information objects.

This supplementary note elaborates the role of trusted-kernel auditing in the overall
optimization-based dependent typing architecture. The recurring theme is that
Pythonic source programs expose many plausible semantic candidates, but the
system must decide which of them should become trusted core obligations, which
should be surfaced as suggested annotations, which should remain as proof
goals, and which should be discarded as too unstable, too expensive, or too
poorly supported. This makes the architecture simultaneously a typing system,
a semantic extraction engine, a theory-integration framework, and an explicit
optimization procedure over candidate dependent information objects.

This supplementary note elaborates the role of trusted-kernel auditing in the overall
optimization-based dependent typing architecture. The recurring theme is that
Pythonic source programs expose many plausible semantic candidates, but the
system must decide which of them should become trusted core obligations, which
should be surfaced as suggested annotations, which should remain as proof
goals, and which should be discarded as too unstable, too expensive, or too
poorly supported. This makes the architecture simultaneously a typing system,
a semantic extraction engine, a theory-integration framework, and an explicit
optimization procedure over candidate dependent information objects.

This supplementary note elaborates the role of trusted-kernel auditing in the overall
optimization-based dependent typing architecture. The recurring theme is that
Pythonic source programs expose many plausible semantic candidates, but the
system must decide which of them should become trusted core obligations, which
should be surfaced as suggested annotations, which should remain as proof
goals, and which should be discarded as too unstable, too expensive, or too
poorly supported. This makes the architecture simultaneously a typing system,
a semantic extraction engine, a theory-integration framework, and an explicit
optimization procedure over candidate dependent information objects.


\section{From first principles: the maximal signature inference problem}

This section states the core problem in the most direct possible way. Suppose a
programmer writes an ordinary Python function or expression and supplies few or
no annotations. What should the system try to infer? The answer cannot merely be
``a nominal type'' or ``a handful of refinements.'' The right target is a
\emph{semantic signature}: a structured object that says as much as possible,
subject to soundness and proof budget, about the value and its relation to the
rest of the program.

\begin{definition}[Semantic signature]
A semantic signature for a term $e$ is a pair
\[
\sigma = \langle \tau, \Phi \rangle
\]
where $\tau$ is a structural or dependent carrier type and $\Phi$ is a finite
family of value-indexed propositions, equations, witnesses, and relational
claims about the denotation of $e$. We write
\[
\sigma(e) \equiv e : \tau \mid \Phi.
\]
Examples include
\[
\py{a : Tensor | issorted(a)},
\qquad
\py{a : Tensor[n] | issorted(a) \wedge distinct(a)},
\]
and, for a function,
\[
\py{sort : (a : Tensor[n]) -> (b : Tensor[n] ** issorted(b) \wedge perm(b,a))}.
\]
\end{definition}

The key point is that the proposition family $\Phi$ is not restricted to scalar
refinements. It may include indexed structure, witness-bearing statements,
protocol facts, heap summaries, and theory-specific facts supplied by a library
module. In particular, the system should be willing to infer not just one fact
such as \(\mathsf{issorted}(a)\), but a whole bundle of mutually supporting
facts whenever they are justified by the code and the active theory bank.

\begin{definition}[Sound inferability]
Let $\mathcal{O}(e)$ denote the observations harvested from the source program,
including guards, branch conditions, control-flow splits, attribute uses,
library calls, tensor operations, protocol constraints, and theorem-pack facts.
A signature $\sigma$ is \emph{soundly inferable} for $e$ when every selected
claim in $\sigma$ is either:
\begin{enumerate}[label=\arabic*.]
  \item derivable in the trusted dependent core,
  \item discharged in an explicitly trusted solver fragment, or
  \item retained as a labeled residual proof obligation.
\end{enumerate}
\end{definition}

The maximal signature inference problem can now be stated cleanly.

\begin{definition}[Maximal automatic signature inference]
Given a term $e$, an observation set $\mathcal{O}(e)$, a candidate-bank
constructor $\mathcal{B}$, a trusted proof boundary $\mathcal{T}$, and a utility
functional $U$, compute
\[
\mathsf{Infer}^\star(e)
= \arg\max_{\sigma \in \mathcal{B}(\mathcal{O}(e))}
U(\sigma)
\]
subject to the sound inferability constraints induced by $\mathcal{T}$.
\end{definition}

This formulation is first-principles in the sense that it does not begin with a
particular syntax of annotations. It begins with the question: what is the most
informative semantic object that can be soundly attached to this term? The
surface syntax of annotations, decorators, contracts, or proof-goal payloads is
secondary. The primary object is the inferred semantic signature.

\subsection{What should count as information?}
A useful signature should say more than ``this is a tensor.'' For a tensor value
$a$, the system may wish to infer members of the following atom families:
\begin{align*}
\mathcal{A}_{\mathrm{shape}}(a) &= \{\mathsf{rank}(a)=r,\ \mathsf{len}(a)=n,\ \mathsf{shape}(a)=\vec{s}\}, \\
\mathcal{A}_{\mathrm{layout}}(a) &= \{\mathsf{dtype}(a)=\delta,\ \mathsf{device}(a)=d,\ \mathsf{contig}(a)\}, \\
\mathcal{A}_{\mathrm{order}}(a) &= \{\mathsf{issorted}(a),\ \mathsf{strictSorted}(a),\ \mathsf{nondecreasing}(a)\}, \\
\mathcal{A}_{\mathrm{set}}(a) &= \{\mathsf{distinct}(a),\ \mathsf{perm}(a,b),\ \mathsf{submultiset}(a,b)\}, \\
\mathcal{A}_{\mathrm{numeric}}(a) &= \{\mathsf{nonnegative}(a),\ \mathsf{bounded}(a,c),\ \mathsf{sum}(a)=s\}, \\
\mathcal{A}_{\mathrm{effect}}(a) &= \{\mathsf{fresh}(a),\ \mathsf{aliases}(a,b),\ \mathsf{requiresGrad}(a)\}.
\end{align*}
For objects and classes, analogous banks include field invariants, ownership
facts, protocol facts, and method-availability facts. For functions, the bank
contains preconditions, postconditions, relational facts between inputs and
outputs, witness-bearing codomains, and theorem-pack-specific obligations.

One immediate implication is that the objective cannot simply maximize the number
of selected atoms. Some atoms are redundant, some are expensive, some activate a
large theory module, some are unstable under refactoring, and some are far more
valuable than others. A signature containing ten trivial restatements of the same
rank fact is worse than one containing rank, sortedness, and a permutation
relation. Therefore the system needs an explicit notion of semantic marginal
value.

\begin{definition}[Semantic utility]
Let $\Phi = \{\phi_1, \ldots, \phi_m\}$ be the selected proposition family of a
candidate signature. The utility of a signature is
\[
U(\sigma)
= \sum_{i=1}^m w_i\,x_i
+ \sum_{1 \le i < j \le m} w_{ij}\,x_i x_j
- \lambda_{\mathrm{ann}} C_{\mathrm{ann}}(\sigma)
- \lambda_{\mathrm{solve}} C_{\mathrm{solve}}(\sigma)
- \lambda_{\mathrm{frag}} C_{\mathrm{frag}}(\sigma)
- \lambda_{\mathrm{unstable}} C_{\mathrm{unstable}}(\sigma)
\]
where $x_i$ indicates selection of proposition $\phi_i$. The pairwise term
captures complementarity: \(\mathsf{issorted}(a)\) may be more valuable when
combined with \(\mathsf{len}(a)=n\) and \(\mathsf{perm}(a,b)\) than in
isolation.
\end{definition}

This is the right place to say explicitly that ``infer as many things as
possible'' does not mean ``infer the largest conjunction in a naive syntactic
sense.'' It means ``infer the highest-utility semantically sound signature.''
The optimization language is essential because the goal is not Boolean
acceptance; it is signature saturation under budgets.

\subsection{The signature lattice and saturation frontier}
Let $\sigma_1 = \langle \tau_1, \Phi_1 \rangle$ and
$\sigma_2 = \langle \tau_2, \Phi_2 \rangle$. We write
\[
\sigma_1 \preceq \sigma_2
\]
when $\tau_2$ is at least as precise as $\tau_1$ and the propositions in
$\Phi_2$ entail those in $\Phi_1$. This gives an information preorder over
signatures. The ideal output would be a greatest sound element above all other
inferable signatures. In practice, however, the greatest element may fail to
exist within the current trusted fragment, or may be too expensive to expose.
The practical target is therefore a \emph{saturation frontier}: the set of
signatures that are undominated with respect to soundness, utility, and budget.

\begin{definition}[Saturation frontier]
The saturation frontier for a term $e$ is
\[
\mathcal{F}(e) = \{\sigma \mid \sigma \text{ is soundly inferable and there is no }
\sigma' \text{ with } U(\sigma') > U(\sigma) \text{ under the same budgets}\}.
\]
\end{definition}

This frontier viewpoint is especially useful for user-facing tooling. One member
of the frontier may correspond to the most compact annotation suggestion. Another
may correspond to the richest proof-goal package. Another may correspond to the
cheapest fully automatic result. A deterministic pipeline can compute or
approximate several frontier points and present them without changing the source
program.

\subsection{Canonical example: inferring \py{a : Tensor | issorted(a)}}
Consider the unannotated function
\begin{lstlisting}[style=deppy,language=Python]
def normalize_sorted(a):
    values, _ = torch.sort(a)
    return values
\end{lstlisting}
A weak system infers only that the return is ``tensor-like.'' A stronger system
should infer at least
\[
(a : \mathsf{Tensor}(n,\delta,d)) \to
(b : \mathsf{Tensor}(n,\delta,d) \mid \mathsf{issorted}(b)).
\]
A still stronger system should infer
\[
(a : \mathsf{Tensor}(n,\delta,d)) \to
(b : \mathsf{Tensor}(n,\delta,d) \mid
\mathsf{issorted}(b) \wedge \mathsf{perm}(b,a)).
\]
An even stronger theory pack could add statements about stability with respect to
equal elements, monotonicity of indices, or preservation of multiplicities. The
point of the optimization formulation is that these increasingly rich signatures
are points on the same inferential spectrum. The system should harvest every
candidate fact, ask which are provable or at least representable, and then solve
for the best package.

\section{Optimization architecture for inferring as many signatures as possible}

We now refine the optimization problem so that it explicitly targets maximal
signature content rather than merely picking a few good predicates. Let
\(\mathcal{P}(e)\) be the property bank for a term or function, let
\(\mathcal{R}(e)\) be the relation bank linking outputs to inputs, let
\(\mathcal{W}(e)\) be the witness-schema bank, and let
\(\mathcal{K}(e)\) be the active library-theory choices. We attach Boolean
variables
\[
 x_p \text{ for } p \in \mathcal{P}(e), \qquad
 y_r \text{ for } r \in \mathcal{R}(e), \qquad
 z_w \text{ for } w \in \mathcal{W}(e), \qquad
 t_k \text{ for } k \in \mathcal{K}(e).
\]
The optimization then becomes a mixed semantic selection problem:
\begin{align*}
\max \quad &
\sum_{p \in \mathcal{P}(e)} \alpha_p x_p
+ \sum_{r \in \mathcal{R}(e)} \beta_r y_r
+ \sum_{w \in \mathcal{W}(e)} \gamma_w z_w
+ \sum_{(u,v) \in \mathcal{I}(e)} \eta_{uv} q_{uv} \\
&- \lambda_1 \sum_p c^{\mathrm{solve}}_p x_p
- \lambda_2 \sum_r c^{\mathrm{solve}}_r y_r
- \lambda_3 \sum_w c^{\mathrm{witness}}_w z_w
- \lambda_4 \sum_k c^{\mathrm{theory}}_k t_k
- \lambda_5 C_{\mathrm{annotation}}(x,y,z)
\end{align*}
subject to compatibility, proof, and dependency constraints. Here the auxiliary
variables $q_{uv}$ linearize useful interactions such as ``sortedness plus
permutation is more valuable than either alone'' or ``shape plus dtype plus
device yields a much better tensor signature than any single fact in isolation.''

The hard constraints have the following form.
\begin{align*}
 x_p &\le \sum_{k \in \mathrm{deps}(p)} t_k, \\
 y_r &\le \sum_{k \in \mathrm{deps}(r)} t_k, \\
 z_w &\le \sum_{k \in \mathrm{deps}(w)} t_k, \\
 x_p &\le \mathsf{Provable}(p) + \mathsf{ResidualAllowed}(p), \\
 y_r &\le \mathsf{Provable}(r) + \mathsf{ResidualAllowed}(r), \\
 z_w &\le \mathsf{Provable}(w) + \mathsf{ResidualAllowed}(w), \\
 q_{uv} &\le u, \qquad q_{uv} \le v, \qquad q_{uv} \ge u + v - 1.
\end{align*}
The first three constraints enforce theory activation; a sortedness property
cannot be selected unless the order theory is active. The next three encode the
trusted boundary. A claim may be selected only if it can be proved in the active
fragment or retained explicitly as a proof obligation. The last line is the
standard linearization for interaction rewards.

\begin{definition}[Maximal signature packing]
For a function $f$, the \emph{maximal signature packing problem} is the problem
of choosing a carrier type, property family, relation family, witness family,
and theory activation vector so as to maximize semantic utility subject to the
hard soundness constraints above.
\end{definition}

The phrase ``packing'' is deliberate. The inferred signature is not a single
logical formula produced by one proof search. It is a packed semantic object made
of structural facts, relational facts, and theory-backed claims whose combined
value exceeds that of the parts. This is why Optimize and MaxSMT are a better
mental model than simple type reconstruction.

\subsection{Why new predicates should keep accumulating}
A practical system cannot freeze its predicate bank. Every new library routine,
every new static bug pattern, and every new theory contribution should enlarge
the space of inferable signatures. We therefore treat the candidate bank as
open-world:
\[
\mathcal{P}_{t+1} = \mathcal{P}_t \cup \Delta\mathcal{P},
\qquad
\mathcal{K}_{t+1} = \mathcal{K}_t \cup \Delta\mathcal{K}.
\]
If a new theory pack adds predicates such as \(\mathsf{issorted}\),
\(\mathsf{perm}\), or \(\mathsf{topkOrdered}\), the optimizer should simply see a
larger feasible region. This is not a reason to rerun an LLM. It is a reason to
rerun the deterministic inference pipeline with a richer candidate basis.

\begin{proposition}[Monotonicity under sound theory accretion]
Suppose a new predicate or theory pack is added without invalidating the
semantics of previously existing packs. Then the feasible set of soundly
inferable signatures can only expand, and the optimal utility value cannot
decrease.
\end{proposition}

\begin{proof}
The old feasible signatures remain feasible because their proof obligations and
dependencies are unchanged. New predicates or theory activations add new
variables and potentially new feasible points. Maximization over a superset of
feasible points cannot yield a worse optimum.
\end{proof}

This proposition is operationally important. It means that the architecture can
continuously grow by adding theory packs, predicate templates, and solver
encodings while keeping the analysis loop deterministic. No new language model
call is required once a theory has been formalized as a candidate generator plus
a trusted or residual-checking interface.

\section{Novel uses of \Zthree{} for signature saturation}

The role of \Zthree{} here should be more ambitious than ``check a few
verification conditions.'' It should act as a programmable engine for signature
saturation, frontier extraction, theory gating, and counterexample analysis.
Below are several uses that go beyond the most basic SMT-backed refinement loop.

\subsection{Lexicographic saturation by semantic tier}
A practical signature should first maximize soundness-critical structure, then
high-value semantic relations, then user-facing compactness. This suggests a
lexicographic objective with tiers such as:
\begin{enumerate}[label=\arabic*.]
  \item maximize definitely proved structural facts,
  \item maximize proved relational and theory facts,
  \item maximize residual-but-explicit proof-goal quality,
  \item minimize annotation cost and solver cost.
\end{enumerate}
\Zthree{} Optimize supports lexicographic priorities directly. This lets the
system prefer, for instance, a signature that proves rank, length, and sortedness
over one that proves only sortedness but leaves basic structure underspecified.

\subsection{Enumerating multiple frontier signatures}
Users and downstream tools often want more than one answer. A human might prefer
a compact signature suggestion; an auditor might prefer the richest proof-goal
package. After obtaining one optimum, the system can add blocker constraints and
enumerate alternative frontier points:
\[
\bigvee_{i : x_i^\star = 1} (x_i = 0)
\ \vee \ 
\bigvee_{i : x_i^\star = 0} (x_i = 1).
\]
Repeated optimization yields a Pareto-style menu of signatures. This is a novel
fit for dependent typing because it treats inference as a frontier-extraction
problem rather than a hunt for one ``principal type.''

\subsection{Unsat-core-guided predicate growth}
When a desired signature fragment cannot be proved, the failure should produce an
unsat core, a countermodel, or both. Instead of merely reporting failure, the
system can solve a secondary hitting-set problem over missing lemmas, missing
theory activations, or missing qualifiers. In effect, \Zthree{} becomes a guide
for deciding which new predicates would most increase the saturation frontier.
This supports continuous theory growth from deterministic artifacts rather than
fresh LLM prompting.

Formally, let $\mathcal{C} = \{C_1, \ldots, C_m\}$ be the family of cores or
counterexample explanations observed across failed proof attempts. Let
$\mathcal{L}$ be the space of candidate repairs. We choose a minimal high-value
repair set by solving:
\[
\min \sum_{\ell \in \mathcal{L}} \kappa_\ell r_\ell
\]
subject to
\[
\forall j \in \{1,\ldots,m\}, \qquad
\sum_{\ell \in \mathrm{covers}(C_j)} r_\ell \ge 1.
\]
This is a weighted hitting-set problem generated by the proof failures
themselves. Once solved, the chosen repairs become new deterministic theory-pack
content.

\subsection{Model-based anti-chain extraction}
When many candidate signatures are mutually incomparable, it is useful to obtain
an anti-chain of high-value interpretations rather than a single dominant point.
\Zthree{} models can be used to detect which properties are simultaneously
realizable and which force tradeoffs. The resulting anti-chain gives the system a
formal explanation for why certain predicates are offered in one suggestion but
not another.

\subsection{Theory gating as a learned but deterministic policy}
Adding every theory pack blindly may be too expensive. We can attach activation
variables $t_k$ to theories and optimize over them as part of the same MaxSMT
problem. Once the cost model is fixed, theory activation is deterministic:
\[
\max \ U(\sigma) - \sum_k \lambda_k t_k.
\]
A new theory pack changes the optimization problem by adding a new $t_k$ and its
associated predicate dependencies, but it does not require rewriting prompts or
re-explaining the source code to an LLM.

\subsection{Countermodel clustering for Python bug classes}
In ordinary Python, many bugs recur in parameterized families: off-by-one index
errors, empty-sequence accesses, wrong tensor-rank assumptions, protocol misses,
heap-field inconsistency, wrong tuple destructuring, and sort-like routines that
return unsorted or shape-mismatched data. \Zthree{} countermodels can be embedded
into a feature space and clustered by explanation shape. The cluster centers can
then induce new deterministic candidate templates such as ``head-like functions
need non-emptiness'' or ``sort-like functions expose sortedness plus
permutation.'' This is a way of discovering theory templates from solver output
without invoking a generative model at analysis time.

\section{Open-world theory packs: keep adding predicates without rerunning the LLM}

A central architectural requirement is that theory growth should be additive,
modular, and deterministic. Once a programmer or researcher has encoded a new
library theory---for example, a PyTorch sorting theory, a broadcasting theory, a
protocol theorem pack, or a heap-ownership fragment---the main inference engine
should simply ingest that pack and recompute the optimization problem. The code
being analyzed does not need to be re-summarized by an LLM. The candidate bank,
proof obligations, and solver encodings are already sufficient.

\begin{definition}[Theory pack]
A theory pack is a tuple
\[
\mathcal{K} = \langle \mathsf{Atoms}, \mathsf{Deps}, \mathsf{Encode}, \mathsf{Trust}, \mathsf{Residualize} \rangle
\]
consisting of atom templates, dependency rules, a lowering function into the
trusted solver fragment when possible, a declaration of the trusted boundary, and
a residualization function for non-lowered obligations.
\end{definition}

This definition is intentionally declarative. Nothing in it requires natural-
language regeneration of the source analysis. Once the pack is installed, the
pipeline changes only by extending $\mathcal{P}(e)$, $\mathcal{R}(e)$,
$\mathcal{W}(e)$, and the dependency graph used by optimization.

\begin{example}[Adding a sortedness theory pack]
Suppose a new theory pack introduces atoms
\[
\mathsf{issorted}(x), \qquad \mathsf{perm}(x,y), \qquad \mathsf{stableSort}(x,y)
\]
and links them to APIs such as \py{torch.sort}, \py{sorted}, and user-defined
merge routines. The pack also supplies a lowering for decidable fragments such as
indexwise monotonicity on bounded vectors and residualization for stronger claims
such as full extensional permutation over symbolic tensors. After installation,
existing Python programs can immediately acquire richer inferred signatures. No
LLM needs to revisit the source file.
\end{example}

\begin{theorem}[Deterministic reanalysis under theory accretion]
Fix a source program, observation extractor, and optimization policy. If theory
packs are added only through declarative pack installation, then the new analysis
result is a deterministic function of the old program and the new pack set.
\end{theorem}

\begin{proof}
The extractor and optimizer are unchanged; only the candidate banks and
encodings are enlarged by installed packs. Therefore the reanalysis result is
computed by the same deterministic pipeline over a larger but well-defined input
basis.
\end{proof}

This property matters because it makes the system scalable as a research and
engineering platform. The language-model component can help write theory packs,
explain diagnostics, or suggest candidate predicates during development, but the
runtime analysis of end-user code does not need to depend on another generation
pass every time the theory library improves.

\subsection{The long-term goal}
The long-term goal is straightforward to state:
\begin{quote}
Given ordinary Python code, infer the richest sound dependent signature available
in the current theory universe, expose the residual proof obligations explicitly,
and improve monotonically as new predicate families and theory packs are added.
\end{quote}
This is the sense in which dependent typing becomes an optimization problem from
first principles. The system is trying to maximize semantic content, not merely
pass or fail a checker. It is trying to infer signatures such as
\py{a : Tensor | issorted(a)} and then to keep going---adding shape, permutation,
device, aliasing, protocol, and heap information whenever the current theories
justify them.



\section{A radically novel judgment form: frontier-valued typing}

Classical typing uses judgments of the form
\[
\Gamma \vdash e : \tau.
\]
Dependent typing enriches this to judgments that mention values, indices, and
proof obligations. The present document argues that, for Python, even this is
still too narrow. The analysis target should be a \emph{frontier} of semantic
signatures. We therefore introduce the judgment
\[
\Gamma ; \mathcal{K} ; B \vdash e \Downarrow \mathcal{F}
\]
where $\mathcal{K}$ is the installed theory-pack set, $B$ is the active proof
and presentation budget, and $\mathcal{F}$ is a saturation frontier of semantic
signatures for $e$.

\begin{definition}[Frontier-valued typing judgment]
The judgment
\[
\Gamma ; \mathcal{K} ; B \vdash e \Downarrow \mathcal{F}
\]
means that, under environment $\Gamma$, theory universe $\mathcal{K}$, and
budget policy $B$, the analysis computes a nonempty set $\mathcal{F}$ of soundly
inferable semantic signatures for $e$ such that no member of $\mathcal{F}$ is
strictly dominated by another feasible signature under the chosen utility
ordering.
\end{definition}

This judgment is radical in a useful way: it formally abandons the expectation
that there is always one best surface type. Instead, it says that the right
semantic output of a practical dependent analysis is a structured tradeoff set.
A programmer may prefer the cheapest automatically proved signature; an auditor
may prefer the richest residual-proof package; a code generator may prefer the
best annotation-compression point. All three may correspond to different members
of the same frontier.

\subsection{Typing as a map into a semantic polytope}
Another way to view the same idea is geometrically. Let the candidate-bank basis
for a term $e$ contain $m$ semantic atoms. Then every binary selection vector
$x \in \{0,1\}^m$ determines a possible semantic package, provided the proof and
dependency constraints are met. Feasible packages therefore form a discrete
polytope inside the hypercube. Typing no longer maps a term to one point; it
maps a term to the exposed face or anti-chain of optima selected by the utility
functional. This suggests the slogan:
\begin{quote}
A practical dependent type checker for Python should be understood as a solver
for a semantic polytope, not merely as a syntax-directed derivation engine.
\end{quote}

\begin{definition}[Semantic signature polytope]
Fix a term $e$ and a candidate basis of size $m$. The \emph{semantic signature
polytope} for $e$ is the set
\[
\mathcal{P}_{\mathrm{sig}}(e) = \{x \in \{0,1\}^m \mid x \text{ satisfies all
proof, compatibility, and theory-dependency constraints}\}.
\]
The frontier $\mathcal{F}(e)$ is the set of exposed optima of
$\mathcal{P}_{\mathrm{sig}}(e)$ under the chosen objective or family of
lexicographic objectives.
\end{definition}

\subsection{Witness amortization}
A particularly novel consequence of the frontier view is that witnesses become
reusable economic objects rather than merely proof artifacts. If the system pays
the proof cost to derive a witness schema for sortedness, permutation, alias
freedom, or protocol conformance, then that witness should be amortizable across
multiple downstream signatures and diagnostics. In a traditional derivation
view, the witness is local to one proof tree. In a frontier view, it is an
asset that can improve multiple frontier points at once. This motivates a new
optimization term:
\[
U_{\mathrm{amortized}}(\sigma) = U(\sigma) + \mu \cdot
\mathsf{ReusePotential}(\mathsf{witnesses}(\sigma)).
\]
Witnesses that unlock many future obligations should therefore be preferred over
witnesses of equal local value but poor reuse.


\section{A formal optimization calculus for dependent signature inference}

This section records the mathematics more explicitly and with less exposition.
The goal is to state a usable formal core for a future implementation.

\subsection{Signature grammar}
We define semantic signatures by the grammar
\begin{align*}
\sigma &::= \langle \Pi, \tau, \Phi, \Omega, W, K \rangle, \\
\Pi &::= \epsilon \mid (x : \tau) , \Pi, \\
\Phi &::= \{\phi_1, \ldots, \phi_m\}, \\
\Omega &::= \{\omega_1, \ldots, \omega_q\}, \\
W &::= \{w_1, \ldots, w_r\}, \\
K &::= \{k_1, \ldots, k_s\}, \\
\phi &::= p(t_1,\ldots,t_n) \mid R(t_1,\ldots,t_n) \mid t = u \mid t \le u \mid \mathsf{Witnessed}(w), \\
\omega &::= \mathsf{Residual}(\psi), \\
w &::= \mathsf{SigmaWitness}(x,\psi) \mid \mathsf{EqWitness}(t,u) \mid \mathsf{TheoryCert}(k,\psi).
\end{align*}
Here $\Pi$ is the parameter telescope, $\tau$ is the carrier type, $\Phi$ is the
selected proposition family, $\Omega$ is the residual obligation family, $W$ is
the witness-schema family, and $K$ is the active theory-pack set.

\subsection{Candidate basis and incidence matrices}
For a term $e$, let the candidate basis be partitioned as
\[
\mathcal{B}(e) = \mathcal{P}(e) \uplus \mathcal{R}(e) \uplus \mathcal{W}(e),
\]
with cardinalities
\[
|\mathcal{P}(e)| = m_p, \qquad |\mathcal{R}(e)| = m_r, \qquad |\mathcal{W}(e)| = m_w.
\]
We use Boolean decision vectors
\[
x \in \{0,1\}^{m_p}, \qquad y \in \{0,1\}^{m_r}, \qquad z \in \{0,1\}^{m_w}, \qquad t \in \{0,1\}^{m_k}
\]
for properties, relations, witnesses, and theories respectively.

Let $A_{\mathrm{dep}} \in \{0,1\}^{(m_p + m_r + m_w) \times m_k}$ be the
dependency matrix with entry $A_{ij}=1$ iff candidate $i$ depends on theory
$k_j$. Let $A_{\mathrm{obs}} \in \{0,1\}^{(m_p + m_r + m_w) \times h}$ record
which observed program facts support which candidates. Let
$A_{\mathrm{frag}} \in \{0,1\}^{(m_p + m_r + m_w) \times c}$ record fragment
requirements such as quantifier freedom, boundedness, linear arithmetic, or
finiteness of the extensional domain.

\subsection{Feasibility constraints}
A selection $(x,y,z,t)$ is feasible only if the following hold.

\paragraph{Theory activation constraints.}
For every candidate index $i$,
\[
s_i \le \sum_{j=1}^{m_k} A_{\mathrm{dep}}[i,j] \cdot t_j,
\]
where $s$ is the concatenation of $(x,y,z)$. If a candidate has no theory
dependencies, the row is treated as vacuously satisfied.

\paragraph{Observation support constraints.}
If the support threshold is $\theta$, then
\[
\forall i, \qquad s_i \le \mathbf{1}\Big[\sum_{j=1}^{h} A_{\mathrm{obs}}[i,j] \ge \theta\Big].
\]
This makes explicit that we do not admit arbitrary semantic fantasies into the
optimization: a candidate must have enough syntactic or theory support.

\paragraph{Fragment admissibility constraints.}
Let $f \in \{0,1\}^c$ indicate which solver fragments are currently allowed.
Then for every candidate index $i$,
\[
s_i \le \prod_{j=1}^{c} f_j^{A_{\mathrm{frag}}[i,j]}.
\]
Operationally, this product is linearized as the family of constraints
\[
\forall j, \qquad s_i \le f_j \quad \text{whenever } A_{\mathrm{frag}}[i,j] = 1.
\]
This is the place where the ``full dependent types, but Z3-decidable where
possible'' requirement enters mathematically. The full language of signatures is
richer than the active automatic fragment, so the optimizer must reason about
fragment boundaries explicitly.

\paragraph{Residualization constraints.}
Each candidate $i$ has a proof status variable
\[
p_i \in \{\Proved, \Refuted, \Unknown\},
\]
encoded concretely by auxiliary Booleans. Selection is allowed only when
\[
s_i \le \mathbf{1}[p_i = \Proved] + r_i,
\]
where $r_i \in \{0,1\}$ is a residualization variable indicating that the
candidate may survive only as an explicit proof obligation. This yields the hard
constraint
\[
s_i \le \pi_i + r_i,
\]
where $\pi_i$ is the Boolean indicator for \Proved.

\subsection{Lexicographic objective}
We now define a lexicographic objective rather than a single weighted sum. Let
\begin{align*}
L_1(x,y,z) &= \sum_{i \in \mathrm{struct}} s_i, \\
L_2(x,y,z) &= \sum_{i \in \mathrm{semantic}} \alpha_i s_i + \sum_{(i,j) \in E_{\mathrm{int}}} \alpha_{ij} q_{ij}, \\
L_3(r) &= \sum_i \rho_i r_i, \\
L_4(x,y,z,t) &= -\sum_i c_i s_i - \sum_j d_j t_j.
\end{align*}
Here $q_{ij}$ are interaction variables satisfying the McCormick constraints
\[
q_{ij} \le s_i, \qquad q_{ij} \le s_j, \qquad q_{ij} \ge s_i + s_j - 1.
\]
The optimization problem is then
\[
\mathrm{lexmax}\ \langle L_1, L_2, L_3, L_4 \rangle
\]
subject to all feasibility constraints.

The intended meaning of the tiers is:
\begin{align*}
L_1 &: \text{maximize proved structural facts}, \\
L_2 &: \text{maximize proved semantic relations and witness value}, \\
L_3 &: \text{maximize high-value residual obligations that should remain visible}, \\
L_4 &: \text{minimize annotation and theory-activation cost.}
\end{align*}
This is more precise than an ordinary weighted objective and aligns directly with
Optimize-style multi-objective solving in \Zthree{}.

\subsection{Frontier extraction}
The frontier is not recovered by one optimization call. Let
\[
s^{\star,1}, s^{\star,2}, \ldots
\]
be a sequence of optimal solutions. After obtaining solution $s^{\star,\ell}$, we
add a blocker clause
\[
\bigvee_{i : s^{\star,\ell}_i = 1} (s_i = 0) \ \vee \ 
\bigvee_{i : s^{\star,\ell}_i = 0} (s_i = 1)
\]
and re-solve under the same lexicographic objective. The extracted frontier is
then the set of nondominated solutions recovered before the objective value drops
below the configured threshold.

\begin{definition}[Algorithmic frontier extraction]
Let $\mathcal{S}_0$ be the feasible set. Recursively define
\begin{align*}
s^{\star,\ell} &= \arg\mathrm{lexmax}_{s \in \mathcal{S}_{\ell-1}} \langle L_1, L_2, L_3, L_4 \rangle, \\
\mathcal{S}_{\ell} &= \mathcal{S}_{\ell-1} \cap \mathsf{Block}(s^{\star,\ell}).
\end{align*}
The extracted frontier after $N$ steps is
\[
\widehat{\mathcal{F}}_N(e) = \{ s^{\star,1}, \ldots, s^{\star,N} \}.
\]
\end{definition}

\subsection{Typing rules as optimization rules}
The frontier-valued judgment can be factored into explicit rules. We write the
most important ones schematically.

\[
\frac{
  \mathcal{O}(e) = \mathsf{Observe}(\Gamma,e)
  \qquad
  \mathcal{B}(e) = \mathsf{Generate}(\mathcal{O}(e), \mathcal{K})
}{
  \Gamma ; \mathcal{K} ; B \vdash e \rightsquigarrow \mathcal{B}(e)
}
\quad [\textsc{Obs-Gen}]
\]

\[
\frac{
  \mathcal{B}(e) \neq \varnothing
  \qquad
  s^\star = \arg\mathrm{lexmax}_{s \in \mathcal{S}(\mathcal{B}(e),B)} \langle L_1,L_2,L_3,L_4 \rangle
}{
  \Gamma ; \mathcal{K} ; B \vdash e \rightsquigarrow s^\star
}
\quad [\textsc{Pack}]
\]

\[
\frac{
  \Gamma ; \mathcal{K} ; B \vdash e \rightsquigarrow s^{\star,1}, \ldots, s^{\star,N}
  \qquad
  \forall \ell\, .\, s^{\star,\ell} \text{ nondominated}
}{
  \Gamma ; \mathcal{K} ; B \vdash e \Downarrow \widehat{\mathcal{F}}_N(e)
}
\quad [\textsc{Frontier}]
\]

\[
\frac{
  \widehat{\mathcal{F}}_N(e) \text{ computed under } \mathcal{K}
  \qquad
  \mathcal{K}' = \mathcal{K} \cup \Delta \mathcal{K}
}{
  \Gamma ; \mathcal{K}' ; B \vdash e \Downarrow \widehat{\mathcal{F}}'_N(e)
}
\quad [\textsc{Accrete}]
\]
with the semantic side condition that the newly installed packs are sound with
respect to their declared trust boundaries.

\subsection{Soundness theorem}
We can now state the basic theorem cleanly.

\begin{theorem}[Soundness of frontier-valued typing]
Assume:
\begin{enumerate}[label=\arabic*.]
  \item the observation extractor is faithful to source semantics,
  \item each theory pack is sound with respect to its declared lowering and
  residualization interface,
  \item the trusted solver fragment proves only semantically valid obligations,
  \item every selected unproved candidate is explicitly residualized.
\end{enumerate}
Then for every signature $\sigma \in \mathcal{F}(e)$ produced by the
frontier-valued judgment, every proposition in $\Phi(\sigma)$ is either
semantically valid or explicitly marked residual in $\Omega(\sigma)$.
\end{theorem}

\begin{proof}[Proof sketch]
Every selected candidate must satisfy the feasibility constraints. Those
constraints require observation support, theory activation, fragment
admissibility, and either proof or residualization. By assumptions (1)--(3), any
candidate marked \Proved is semantically sound. By assumption (4), any selected
candidate that is not proved appears in the residual set. Therefore no signature
contains silently unsound selected facts.
\end{proof}

\subsection{Monotonicity under theory extension}
We can also state more formally what theory accretion buys us.

\begin{theorem}[Frontier monotonicity under sound theory extension]
Let $\mathcal{K} \subseteq \mathcal{K}'$ and suppose every pack in
$\mathcal{K}' \setminus \mathcal{K}$ is sound and only adds new candidate atoms,
relations, witnesses, and dependencies without invalidating previous ones. Then
for every feasible signature under $\mathcal{K}$ there exists a feasible
signature under $\mathcal{K}'$ with utility at least as large. Consequently, the
best lexicographic objective value under $\mathcal{K}'$ is at least as good as
under $\mathcal{K}$.
\end{theorem}

\begin{proof}
Reuse the old decision vector and set all new candidate and theory-activation
variables to zero. The old solution remains feasible because previously valid
constraints are unchanged. Since the new feasible region strictly contains the old
one, optimization over the larger region cannot reduce the optimum.
\end{proof}

\subsection{Unsat-core-driven repair optimization}
Finally, when proof fails, we define repair as an optimization problem in its own
right. Let $\mathcal{C} = \{C_1, \ldots, C_n\}$ be the observed family of unsat
cores or counterexample explanation sets, and let $\mathcal{L} = \{\ell_1,
\ldots, \ell_g\}$ be the candidate repair set of new lemmas, qualifiers, witness
schemas, or theory activations. Define cover matrix
$H \in \{0,1\}^{n \times g}$ by
\[
H[i,j] = 1 \iff \ell_j \text{ repairs explanation } C_i.
\]
We then solve
\[
\min_{r \in \{0,1\}^g} \sum_{j=1}^{g} \kappa_j r_j
\]
subject to
\[
\forall i \in \{1,\ldots,n\}, \qquad \sum_{j=1}^{g} H[i,j] r_j \ge 1.
\]
The chosen repair set induces
\[
\Delta \mathcal{K} = \mathsf{Install}(r),
\qquad
\Delta \mathcal{P} = \mathsf{NewAtoms}(r),
\]
which then feeds the next frontier computation. This closes the mathematical loop
between checking, optimization, and theory growth.


\section{Two-source inference: explicit LLM semantics plus unannotated user code}

The system this monograph advocates must serve two extreme but complementary
modes of use.
\begin{enumerate}[label=\arabic*.]
  \item An LLM or expert user may provide an explicit semantic contract,
  candidate proof objects, theorem schemas, or highly detailed dependent type
  information.
  \item An ordinary user may provide almost no annotations at all, expecting the
  system to infer as much semantic structure as possible from ordinary Python,
  PyTorch, or library-heavy code.
\end{enumerate}
The mathematics should not bifurcate into two unrelated pipelines. Both modes
should reduce to one optimization problem over a shared candidate basis.

\subsection{Source-partitioned observation model}
Let $e$ be the source program and let $m$ denote optional machine-supplied or
LLM-supplied semantic metadata. We partition observations into three families:
\begin{align*}
\mathcal{O}_{\mathrm{code}}(e) &:= \text{facts harvested from source syntax,
guards, control flow, calls, and dataflow}, \\
\mathcal{O}_{\mathrm{llm}}(m) &:= \text{facts supplied explicitly by the LLM or
expert user}, \\
\mathcal{O}_{\mathrm{theory}}(e,\mathcal{K}) &:= \text{facts induced by installed
theory packs over observed library operations}.
\end{align*}
The total observation universe is
\[
\mathcal{O}_{\mathrm{tot}}(e,m,\mathcal{K})
= \mathcal{O}_{\mathrm{code}}(e)
\cup \mathcal{O}_{\mathrm{llm}}(m)
\cup \mathcal{O}_{\mathrm{theory}}(e,\mathcal{K}).
\]
Candidate generation then becomes
\[
\mathcal{B}(e,m,\mathcal{K})
= \mathsf{Generate}(\mathcal{O}_{\mathrm{tot}}(e,m,\mathcal{K})).
\]

This already captures the desired unification: explicit LLM-provided facts and
latent code-mined facts are not separate species of object. They differ only in
provenance, confidence, trust pathway, and residualization policy.

\subsection{Source-weighted utility}
We now refine the utility function so that explicit LLM claims can be rewarded
for semantic richness without being trusted blindly. Let every candidate $i$ have
source label
\[
\mathsf{src}(i) \in \{\mathrm{code}, \mathrm{llm}, \mathrm{theory}\}
\]
and let $\lambda_{\mathrm{src}}$ be a source-specific confidence multiplier. The
utility of a selected package is then
\[
U(\sigma)
= \sum_i \lambda_{\mathsf{src}(i)} \alpha_i s_i
+ \sum_{(i,j) \in E_{\mathrm{int}}} \beta_{ij} q_{ij}
- \sum_i c_i s_i
- \sum_k d_k t_k
- \sum_i \chi_i u_i.
\]
The last term penalizes unchecked uptake of externally supplied facts, where
$u_i$ is an uncertainty surcharge variable. In particular, an LLM-supplied fact
may have high semantic value but should still incur a penalty unless it is
checked, lowered, or explicitly residualized.

\subsection{Consistency constraints between explicit and latent facts}
Suppose the LLM proposes a candidate family $\mathcal{L} \subseteq
\mathcal{B}(e,m,\mathcal{K})$ containing, for example, a precise dependent
signature, a proof sketch of sortedness, or a permutation certificate for a
sorting routine. Let the latent code-mined family be $\mathcal{C}$. The optimizer
must enforce cross-source consistency. We model this with a consistency matrix
$M \in \{-1,0,1\}^{|\mathcal{L}| \times |\mathcal{C}|}$, where
\begin{align*}
M[i,j] = 1 &\iff \ell_i \text{ and } c_j \text{ reinforce each other}, \\
M[i,j] = -1 &\iff \ell_i \text{ and } c_j \text{ are inconsistent}, \\
M[i,j] = 0 &\iff \text{no direct interaction is known.}
\end{align*}
We then enforce
\[
\forall (i,j) \text{ with } M[i,j] = -1, \qquad \ell_i + c_j \le 1,
\]
and reward reinforcement through interaction variables
\[
q_{ij} \le \ell_i, \qquad q_{ij} \le c_j, \qquad q_{ij} \ge \ell_i + c_j - 1.
\]
This yields a principled way to combine a strong LLM-supplied theorem sketch with
independent evidence from the user's actual code. The result is neither blind
acceptance nor blind distrust.

\subsection{Semantic proof objects as candidates}
If the LLM supplies proof terms, proof sketches, certificates, or witness
schemas, these should enter the same optimization space as ordinary predicates.
Let
\[
\mathcal{W}_{\mathrm{llm}}(m) = \{w_1, \ldots, w_r\}
\]
be the family of proposed semantic proof objects. Each witness variable $z_j$
then satisfies
\[
z_j \le \pi_j + r_j,
\]
where $\pi_j$ indicates successful checking in the trusted kernel or solver
fragment, and $r_j$ indicates explicit residualization. This makes witness uptake
obey the same soundness discipline as ordinary predicates while still allowing
semantically rich LLM output to improve the final signature frontier.

\subsection{Joint objective for annotated and unannotated modes}
The two user modes can now be seen as parameter settings of one system.
\begin{enumerate}[label=\arabic*.]
  \item In the \emph{unannotated mode}, $\mathcal{O}_{\mathrm{llm}}(m) = \varnothing$,
  so optimization ranges only over code-mined and theory-induced candidates.
  \item In the \emph{LLM-explicit mode}, $\mathcal{O}_{\mathrm{llm}}(m)$ may be rich,
  contributing detailed dependent signatures, semantic theorems, and witness
  schemas.
  \item In the \emph{mixed mode}, the optimizer fuses both sources and chooses
  the strongest coherent frontier under the trust and cost budget.
\end{enumerate}
Thus ``unannotated'' and ``fully explicit'' are not separate systems; they are
points on one optimization continuum.

\subsection{Formal soundness principle for LLM-assisted inference}
We can state the critical theorem directly.

\begin{theorem}[Conservative soundness under LLM assistance]
Assume that every LLM-supplied candidate is admitted into the final signature only
if it is either:
\begin{enumerate}[label=\arabic*.]
  \item checked successfully in the trusted kernel,
  \item checked successfully in an approved solver fragment, or
  \item surfaced as an explicit residual obligation.
\end{enumerate}
Then the addition of LLM-supplied semantic information cannot make the final
system silently unsound.
\end{theorem}

\begin{proof}
The proof is identical in shape to the basic soundness theorem, except that the
candidate basis now includes an additional provenance class. Since provenance
alone does not bypass proof or residualization constraints, every selected
LLM-supplied fact is either checked or made explicit as residual. Therefore no
silent unsoundness is introduced by the additional source.
\end{proof}

\subsection{Canonical dual-mode example}
A sorting routine is the canonical example because it exposes both regimes.

\paragraph{LLM-explicit regime.}
The LLM proposes
\[
(xs : \mathsf{Tensor}(n,\delta,d)) \to
(result : \mathsf{Tensor}(n,\delta,d) \mid
\mathsf{sorted}(result) \wedge \mathsf{perm}(result,xs)).
\]
It may also propose a witness that the returned tensor is a stable ordering of
the input. These become candidate atoms and witness schemas with high semantic
value but explicit trust requirements.

\paragraph{Unannotated-user regime.}
The user writes only
\begin{lstlisting}[style=deppy,language=Python]
def sort_values(xs):
    values, _ = torch.sort(xs)
    return values
\end{lstlisting}
The system mines the latent code facts, applies the installed sorting theory,
and recovers a closely related signature even when no LLM metadata is supplied.

\paragraph{Mixed regime.}
The LLM may supply a stronger theorem---for example, stability or a domain-
specific permutation notion---while the user's actual code still determines
whether the theorem is coherent with the implementation. Both sources are fused
by one optimization problem and one trust discipline.


\section{A dual-surface calculus: F$^\star$-style explicit users and ordinary Python users}

The intended system should not merely accept two sources of evidence. It should
support two different \emph{user surfaces} that elaborate into one common
semantic core.

\subsection{Two source languages, one core}
Let
\[
\mathcal{S}_{\mathrm{proof}}
\qquad\text{and}\qquad
\mathcal{S}_{\mathrm{py}}
\]
be two source languages.

\paragraph{Proof-oriented surface.}
$\mathcal{S}_{\mathrm{proof}}$ is the surface for users---including LLMs---who
want to use the system as if it were an F$^\star$-like semantic programming and
proof environment. Its terms may contain explicit dependent signatures, witness
terms, theorem declarations, refinement predicates, and proof-carrying library
contracts. Abstractly:
\begin{align*}
M &::= \mathsf{def}\ f : \Pi \to \tau\ \{e\}
\mid \mathsf{theorem}\ h : \phi\ \{p\}
\mid \mathsf{lemma}\ \ell : \phi
\mid \mathsf{theory}\ k\{\Delta\}, \\
\Pi &::= \epsilon \mid (x : A)\,\Pi, \\
A &::= \Type \mid \Pi x : A . B \mid \Sigma x : A . B \mid \mathsf{Eq}(A,t,u)
\mid \{x : A \mid \phi\} \mid \mathsf{Tensor}(\vec{s},\delta,d) \mid \cdots.
\end{align*}
This surface is meant for users who want to be maximally explicit.

\paragraph{Pythonic surface.}
$\mathcal{S}_{\mathrm{py}}$ is the surface for regular users writing ordinary
Python or PyTorch code with any annotation level from zero to full. Abstractly:
\begin{align*}
P &::= \mathsf{def}\ f(x_1,\ldots,x_n) :?\ \{b\}
\mid \mathsf{class}\ C\{\Delta\}
\mid \mathsf{import}\ L
\mid \mathsf{stmt}, \\
:? &::= \epsilon \mid \to A \mid \mathsf{sig}(\Pi \to A),
\end{align*}
where $:?$ indicates that the user may provide no annotation, a shallow Python
annotation, or a rich dependent signature.

\paragraph{Common semantic core.}
Both surfaces elaborate into a common dependent core
\[
\mathcal{C}_{\mathrm{dep}}
\]
whose objects are semantic signatures, residual obligations, witness schemas,
and theory activations.

\subsection{Two elaboration judgments}
The explicit surface is elaborated by a checking-oriented judgment
\[
\Gamma ; \mathcal{K} \vdash_{\mathrm{proof}} M \leadsto c : \sigma
\]
where $c \in \mathcal{C}_{\mathrm{dep}}$ is a fully elaborated semantic object.
Because the user already supplied much of the semantic structure, this judgment
is closer to traditional dependent elaboration.

The Pythonic surface is elaborated by an optimization-oriented judgment
\[
\Gamma ; \mathcal{K} ; B \vdash_{\mathrm{py}} P \leadsto \widehat{\mathcal{F}}(P)
\]
which returns a frontier of candidate semantic signatures. The elaboration is not
just type reconstruction; it is observation, candidate generation, optimization,
and residualization.

The point of the architecture is that both judgments target the same semantic
objects. The difference is not the destination, but the amount of information
that must be synthesized on the way there.

\subsection{The explicit surface as a degenerate optimization problem}
A useful formal simplification is to treat the proof-oriented surface as a
special case of the optimization problem. Let the user-supplied semantic package
induce a fixed candidate subset
\[
S_{\mathrm{user}} \subseteq \mathcal{B}(M,\mathcal{K}).
\]
Then the explicit mode can be modeled by the constrained optimization problem
\[
\max U(s)
\]
subject to
\[
\forall i \in S_{\mathrm{user}}, \quad s_i = 1,
\qquad
\forall j \in S_{\mathrm{forbidden}}, \quad s_j = 0,
\]
plus the usual proof, dependency, and residualization constraints. Thus the
proof-oriented user is not leaving the optimization framework; they are injecting
hard constraints into it.

This yields the key unification principle:
\begin{quote}
F$^\star$-style users and ordinary Python users interact with the same semantic
optimizer. The former pin many variables explicitly; the latter leave more of the
space free for optimization to discover.
\end{quote}

\subsection{Annotation level as a continuous control variable}
To make the spectrum precise, let $a \in [0,1]$ denote the annotation-explicitness
level of a module or function.
\begin{align*}
a = 0 &\quad\text{means fully unannotated or nearly so}, \\
a = 1 &\quad\text{means fully explicit proof-oriented programming.}
\end{align*}
We can then make the optimization itself depend on $a$. For example, let
\[
\lambda_{\mathrm{user}}(a) = \lambda_0 + \eta a
\]
control how strongly user-supplied explicit facts are preferred, while
\[
\lambda_{\mathrm{discover}}(a) = \mu_0 + \mu_1 (1-a)
\]
controls how aggressively the optimizer rewards newly discovered latent facts.
The objective becomes
\[
U_a(s)
= \lambda_{\mathrm{user}}(a) \cdot U_{\mathrm{explicit}}(s)
+ \lambda_{\mathrm{discover}}(a) \cdot U_{\mathrm{latent}}(s)
- C(s).
\]
At $a \approx 1$, the system behaves more like an explicit proof assistant. At
$a \approx 0$, it behaves more like a semantics-seeking static analyzer. In the
middle, it interpolates smoothly between the two.

\subsection{Conservative extension theorem}
We can now state the formal relation between the two surfaces.

\begin{theorem}[Conservative extension of the proof-oriented surface]
Suppose a module $M \in \mathcal{S}_{\mathrm{proof}}$ elaborates successfully to a
semantic signature $\sigma$ under the proof-oriented judgment. Then there exists a
presentation of $M$ inside the general optimization framework such that the
resulting frontier contains a distinguished point $\sigma^\star$ equivalent to
$\sigma$, and no additional optimization freedom is required to recover it.
\end{theorem}

\begin{proof}[Proof sketch]
Encode all user-supplied semantic commitments as hard equalities over the
selection variables and witness variables. Because the proof-oriented elaboration
already determines the semantic structure, the feasible set of the induced
optimization problem contains a point corresponding to that structure. The proof
obligations and residualization behavior are identical, so the recovered point is
semantically equivalent.
\end{proof}

\begin{corollary}
Adding the Pythonic optimization machinery does not weaken the system for users
who want to work in a proof-first, F$^\star$-like style. It strictly generalizes
that mode.
\end{corollary}

\subsection{Lower-bound theorem for ordinary Python users}
The system should also provide a guarantee in the opposite direction: regular
users should get something meaningful even when they supply little or nothing.

\begin{theorem}[Latent semantic lower bound]
Assume that the observation extractor is complete for a finite family of local
patterns $\Lambda$ (guards, index uses, attribute requirements, theory-triggering
calls, and destructuring patterns), and that each pattern in $\Lambda$ maps to at
least one candidate semantic atom. Then every Python program containing at least
one matched pattern has a nontrivial feasible signature whose information content
is strictly greater than that of the trivial top-type signature.
\end{theorem}

\begin{proof}
By assumption, every matched pattern contributes at least one supported candidate
atom. Since the candidate has support and is admitted by the finite pattern basis,
it yields a feasible nontrivial selection unless contradicted by hard proof or
compatibility constraints. Therefore the trivial all-top signature is no longer
the unique feasible point.
\end{proof}

This theorem is modest but important: it says ordinary Python users are not just
using a degraded proof assistant. They are users of a distinct but formally
integrated mode whose analysis still yields semantic lower bounds.

\subsection{Bimodal example schema}
The same library routine should be expressible in both surfaces.

\paragraph{Proof-oriented presentation.}
A user or LLM may write, in effect,
\[
\mathsf{sort} : (xs : \mathsf{Tensor}(n,\delta,d)) \to
(result : \mathsf{Tensor}(n,\delta,d) \mid
\mathsf{sorted}(result) \wedge \mathsf{perm}(result,xs)).
\]
They may further attach a semantic witness certifying stability or monotonicity.

\paragraph{Ordinary Python presentation.}
A regular user writes only
\begin{lstlisting}[style=deppy,language=Python]
def sort_values(xs):
    values, _ = torch.sort(xs)
    return values
\end{lstlisting}
The optimizer should recover a frontier containing a point close to the explicit
signature above, modulo the currently installed theories and proof budget.

\paragraph{Design consequence.}
The real system target is therefore not ``an analyzer with optional LLM hints.''
It is a \emph{bimodal dependent environment}: one core, two surfaces, one trust
discipline, and one optimization problem over semantic information.

\section{A clean-slate implementation agenda for the eventual rewrite}

Because the user explicitly allows a clean-slate restart of infrastructure, it is
worth stating what the implementation should ultimately optimize for if rebuilt
from scratch.

\subsection{Architectural principle 1: semantic signatures are the primary IR}
A future implementation should not treat contracts, minimal annotations,
diagnostics, proof goals, and theory facts as separate post hoc views glued
together late in the pipeline. Instead, it should construct a semantic-signature
intermediate representation as early as possible. Every later artifact should be
a rendering or projection of that IR.

\subsection{Architectural principle 2: theory packs should be data, not hard-coded branches}
Library support should be declarative. A theory pack should specify atom
templates, dependency rules, trust boundaries, residualization rules, and
optional solver lowerings. The runtime pipeline should simply load installed
packs and solve the enlarged optimization problem. This makes the system truly
open-world.

\subsection{Architectural principle 3: one optimizer, many renderings}
The same packed signature should feed at least five outward-facing products:
\begin{enumerate}[label=\arabic*.]
  \item a rich semantic signature view for researchers and auditors,
  \item a minimal annotation suggestion for programmers,
  \item explicit proof goals for theorem-prover or human follow-up,
  \item static diagnostics for concrete bug reports,
  \item counterexample and frontier traces for improving theory packs.
\end{enumerate}
A future rewrite should therefore have exactly one semantic optimizer and many
thin renderers, rather than many analyses that each rediscover overlapping
facts.

\subsection{Architectural principle 4: full dependent ambition, explicit automation boundary}
The implementation should aim at full dependent expressivity internally, while
also making explicit what falls inside the trusted automatic fragment. The goal
is not to collapse everything into SMT. The goal is to use SMT aggressively for
selection, packing, and decidable proof discharge while preserving residual
obligations as first-class objects.

\subsection{Architectural principle 5: unannotated-to-fully-annotated is a spectrum}
A good system should let users live anywhere on the spectrum between zero
annotations and theorem-prover-level explicit signatures. The monograph's
mathematics already implies this: an annotation is just a user-supplied point,
constraint, or preference inside the larger semantic-signature search space. The
future implementation should therefore treat the phrases ``unannotated'' and
``fully annotated'' not as different modes, but as different priors over the
same underlying inference engine.

\section{Capstone example: unannotated PyTorch inference beyond TensorGuard}

We end with the strongest current example because it captures the actual target
of this project more honestly than any slogan. The goal is not merely to retain
TensorGuard-style shape checking, valuable though that is. The goal is to infer
\emph{semantic} signatures for ordinary unannotated PyTorch code and then use
those inferred signatures to find bugs statically that shape-only reasoning would
miss.

Consider the following completely unannotated code.
\begin{lstlisting}[style=deppy,language=Python]
def infer_sorted_values(xs):
    values, _indices = torch.sort(xs, dim=-1)
    return values


def buggy_sort_or_indices(xs, take_indices):
    values, indices = torch.sort(xs, dim=-1)
    result = values
    if take_indices:
        result = indices
    return result
\end{lstlisting}

A weak analysis says only that both functions manipulate tensors. A shape-centric
analysis in the style of TensorGuard may reason about dimensional consistency,
but it does not naturally distinguish the semantic role of the first projection
of \py{torch.sort} from the semantic role of the second projection. By contrast,
our optimization-based dependent view treats \py{torch.sort} as a theory-bearing
library primitive that generates multiple candidate facts:
\begin{enumerate}[label=\arabic*.]
  \item the first projection is a value tensor,
  \item the first projection is sorted,
  \item the first projection is a permutation of the input,
  \item the second projection is an index tensor,
  \item mixing the two should not be silently accepted.
\end{enumerate}

In the current implementation, the strongest automatically selected signature for
\py{infer\_sorted\_values} is approximately
\[
\mathsf{infer\_sorted\_values} :
(xs : \mathsf{Tensor}) \to
(result : \mathsf{SortedTensor} \mid \mathsf{sorted}(result) \wedge \mathsf{perm}(result, xs)).
\]
This matters for two reasons. First, the signature is inferred from unannotated
source. Second, it is already stronger than what TensorGuard-style shape
assertions usually provide, because it captures a semantic relation between the
input and the output rather than only a dimensional discipline.

Now consider \py{buggy\_sort\_or\_indices}. The analysis infers that
\py{values} has type \py{SortedTensor} while \py{indices} has type
\py{TensorIndex}. The assignment
\begin{lstlisting}[style=deppy,language=Python]
result = indices
\end{lstlisting}
therefore becomes a static semantic mismatch, because \py{result} was already
established as a sorted value tensor. The system reports an
\py{assignment-mismatch} without requiring any user annotation. Intuitively, the
program has confused two different semantic products of sorting: the sorted data
and the permutation witness.

This is exactly the sort of example that motivates the full optimization view.
The system is not merely checking one provided contract. It is harvesting a
library theory, generating candidate semantic facts, selecting the highest-value
sound subset, and then using those inferred facts to reject a bug. The winning
point is not that this particular example uses \py{torch.sort}. The winning
point is that once the theory pack for sorting exists, the same deterministic
pipeline can extend to other operators---\py{topk}, \py{argsort}, \py{gather},
shape-changing views, broadcasting operators, attention primitives, and domain-
specific kernels---without rerunning an LLM over the program text.

If this project succeeds, then this example will stop looking special. It will
instead look like the first member of a large family of unannotated Python and
PyTorch programs for which the system can infer not just ``tensorness'' or
``shape compatibility,'' but rich dependent signatures and semantic bug reports.
That is the destination: ordinary code, maximal inferred meaning, solver-backed
selection, continuously growing theory packs, and static bug finding that goes
well beyond TensorGuard.

\end{document}
